% SEGMENT 2 of 2: Pages 46-85
% Catégories tannakiennes by P. Deligne
% Typeset from pages 46-85 (original publication pages 156-195)

\documentclass[12pt,a4paper]{article}
\usepackage[T1]{fontenc}
\usepackage[utf8]{inputenc}
\usepackage[french]{babel}
\usepackage{amsmath, amssymb, amsthm, amscd}
\usepackage{mathtools}
\usepackage{enumitem}
\usepackage{graphicx}
\usepackage{xcolor}
\usepackage[margin=2.5cm]{geometry}
\usepackage{hyperref}
\usepackage{tikz-cd}
\usepackage{mathrsfs}

% Theorem environments
\theoremstyle{plain}
\newtheorem{theorem}{Théorème}[section]
\newtheorem{lemma}[theorem]{Lemme}
\newtheorem{proposition}[theorem]{Proposition}
\newtheorem{corollary}[theorem]{Corollaire}

\theoremstyle{definition}
\newtheorem{definition}[theorem]{Définition}
\newtheorem{remark}[theorem]{Remarque}
\newtheorem{example}[theorem]{Exemple}

% Math operators
\DeclareMathOperator{\Spec}{Spec}
\DeclareMathOperator{\Sp}{Sp}
\DeclareMathOperator{\Hom}{Hom}
\DeclareMathOperator{\End}{End}
\DeclareMathOperator{\Aut}{Aut}
\DeclareMathOperator{\Isom}{Isom}
\DeclareMathOperator{\Rep}{Rep}
\DeclareMathOperator{\Id}{Id}
\DeclareMathOperator{\Vect}{Vect}
\DeclareMathOperator{\ev}{ev}
\DeclareMathOperator{\tr}{Tr}
\DeclareMathOperator{\Sym}{Sym}
\DeclareMathOperator{\Ind}{Ind}
\DeclareMathOperator{\coker}{coker}
\DeclareMathOperator{\im}{Im}
\DeclareMathOperator{\rk}{rang}
\DeclareMathOperator{\coim}{coim}
\DeclareMathOperator{\diag}{diag}
\DeclareMathOperator{\ptf}{ptf}
\DeclareMathOperator{\pr}{pr}
\DeclareMathOperator{\inj}{inj}
\DeclareMathOperator{\Der}{Der}
\DeclareMathOperator{\GL}{GL}
\DeclareMathOperator{\Ob}{Ob}
\DeclareMathOperator{\Tr}{Tr}
\newcommand{\Ens}{\mathrm{Ens}}
\newcommand{\coh}{\mathrm{coh}}
\newcommand{\eed}{\text{erd}}

\newcommand{\ot}{\otimes}
\newcommand{\fbin}{\times_{\text{fin}}}
\newcommand{\dirlim}{\varinjlim}
\newcommand{\invlim}{\varprojlim}

\errorcontextlines=100
\begin{document}
\nonstopmode
\setcounter{section}{5}
\setcounter{theorem}{4}

\subsection{}

Pour $\omega_1(X \otimes Y) \xrightarrow{\sim} \omega_1(X) \otimes \omega_1(Y)$. Par définition d'un $\otimes$-foncteur, l'isomorphisme de foncteurs composés rendant commutatif le bord du diagramme est le même qu'en (4), une fois identifié $\omega(Y \otimes X)$ à $\omega(X \otimes Y)$ par la commutativité de $\otimes$ dans $\mathcal{A}$. Appliquant $L(\, , \,)$ à ce diagramme, on obtient la commutativité de (6.3.2). L'associativité s'obtient de même. Si $\{1\}$ est la sous-catégorie de $\mathcal{A}$ réduite à l'objet unité et à sa flèche identique, on a, par définition des $\otimes$-foncteur $\omega_j(1) \simeq B$ et $L_j(\omega_1|\{1\}, \omega_2|\{1\}) = B$. L'unité $B \to L_j(\omega_1, \omega_2)$ de $L_j(\omega_1, \omega_2)$ est définie par 4.10(ii).

\subsection{}

Pour $\mathcal{A}$ une catégorie à produit tensoriel vérifiant (2.1.1) et $\omega_1, \omega_2$ deux $\otimes$-foncteurs de $\mathcal{A}$ dans les modules sur un schéma $S$, nous noterons $\underline{\Hom}^\otimes_S(\omega_1, \omega_2)$ (resp.~$\underline{\Isom}^\otimes_S(\omega_1, \omega_2)$) le foncteur sur $(\mathrm{Sch}/S)$ qui à $u : T \to S$ attache l'ensemble des morphismes (resp.~isomorphismes) de $\otimes$-foncteurs $u^*\omega_1 \to u^*\omega_2$ (2.7). Si $\omega_1$ et $\omega_2$ sont à valeurs dans les modules localement libres de rang fini, ces foncteurs sont représentables par un schéma affine sur $S$. D'après 2.7, si $\mathcal{A}$ vérifie (2.1.1) et (2.1.2), cette condition est vérifiée.

Comme en 1.11, pour $\omega_i$ à valeurs dans les modules sur $S$, on pose
\[
\underline{\Hom}^k_k(\omega_2, \omega_1) := \underline{\Hom}^k_{S \times_k S}(\pr_1^*\omega_2, \pr_2^*\omega_1).
\]
De même pour $\underline{\Isom}$. Pour $\omega_1 = \omega_2$, on écrit $\underline{\End}(\omega)$ pour $\underline{\Hom}(\omega, \omega)$ et $\underline{\Aut}(\omega)$ pour $\underline{\Isom}(\omega, \omega)$.

\begin{proposition}[6.6]
Soient $\mathcal{A}$ une catégorie à produit tensoriel vérifiant (2.1.1) et $\omega_1, \omega_2$ deux $\otimes$-foncteurs de $\mathcal{A}$ dans $(B\text{-}\mod)_{\ptf}$. Posons $S := \Spec(B)$.

Le schéma $\Spec L_S(\omega_1, \omega_2)$ sur $S$ représente le foncteur $\underline{\Hom}^\otimes_S(\omega_2, \omega_1)$.
\end{proposition}

\begin{proof}
Soient $T = \Spec(C)$ un schéma affine sur $S$, $u : T \to S$. Par définition (4.7) un morphisme $f$ de $B$-modules de $L_S(\omega_1, \omega_2)$ dans $C$ s'identifie à un système fonctoriel de morphismes de $B$-modules
\[
f_X : \omega_2(X) \to \omega_1(X) \otimes C.
\]
La donnée des $f_X$ équivaut à celle de morphismes $C$-linéaires
\[
f'_X : \omega_2(X) \otimes C \to \omega_1(X) \otimes C,
\]
fonctoriels en $X$, i.e.~à un morphisme $f'$ de foncteurs de $u^*\omega_2$ dans $u^*\omega_1$.

On laisse au lecteur le soin de vérifier que pour que $f$ soit un morphisme d'algèbres, il faut et il suffit que $f'$ soit un morphisme de $\otimes$-foncteurs.
\end{proof}

\subsection{}

Soient $k$ un anneau commutatif et $B_1$ et $B_2$ deux $k$-algèbres commutatives. Soient $\mathcal{A}$ une catégorie à produit tensoriel vérifiant (2.1.1) et $L_j$ ($j=1,2$) un $\otimes$-foncteur de $\mathcal{A}$ dans $(B_j\text{-}\mod)_{\ptf}$. Par extension des scalaires, $\omega_1$ (resp.~$\omega_2$) définit un $\otimes$-foncteur de $\mathcal{A}$ dans $(B_1 \otimes_k B_2\text{-}\mod)_{\ptf}$; on le note $\omega_1 \otimes 1$ (resp.~$1 \otimes \omega_2$). On a
\[
L_k(\omega_1, \omega_2) = L_{B_1 \otimes_k B_2}(\omega_1 \otimes 1, 1 \otimes \omega_2).
\]

D'après 6.6, $\Spec_k L_k(\omega_1, \omega_2)$ représente le foncteur $\underline{\Hom}^k_k(\omega_2, \omega_1)$. D'après 6.5, ce dernier coïncide avec $\underline{\Isom}^k_k(\omega_1, \omega_2)$ si $\mathcal{A}$ vérifie (2.1.2).

Pour trois $B_j$ et $\omega_j$, la composition des morphismes
\[
\underline{\Hom}^k_3(\omega_3, \omega_2) \times \underline{\Hom}^k_2(\omega_2, \omega_1) \to \underline{\Hom}^k_2(\omega_3, \omega_1)
\]
correspond à un morphisme de $k$-algèbres
\[
c : L_k(\omega_1, \omega_3) \to L_k(\omega_1, \omega_2) \otimes_k L_k(\omega_2, \omega_3).
\]

Par définition, le morphisme de $L_k(\omega_1, \omega_2) \otimes_k L_k(\omega_2, \omega_3)$-modules déduit par extension des scalaires (par $c$) de
\[
\omega_3(X) \otimes_k L_k(\omega_1, \omega_3) \to \omega_1(X) \otimes_k L_k(\omega_1, \omega_3)
\]
est le composé des morphismes déduits par extension des scalaires de
\[
\omega_3(X) \otimes_k L_k(\omega_2, \omega_3) \to \omega_2(X) \otimes_k L_k(\omega_2, \omega_3)
\]
et
\[
\omega_2(X) \otimes_k L_k(\omega_1, \omega_2) \to \omega_1(X) \otimes_k L_k(\omega_1, \omega_2).
\]

Ceci revient à la commutativité
\[
\begin{tikzcd}
\omega_3(X) \arrow[r] \arrow[d] & \omega_1(X) \otimes_k L_k(\omega_1, \omega_2) \arrow[d] \\
\omega_2(X) \otimes_k L_k(\omega_2, \omega_3) \arrow[r] & \omega_1(X) \otimes_k L_k(\omega_1, \omega_2) \otimes_k L_k(\omega_2, \omega_3)
\end{tikzcd}
\]
et $c$ est donc (4.7.4).

\subsection{Preuve de 1.12(ii) pour $S$ affine}

Soient $\mathcal{T}$ une catégorie tensorielle sur $k$, $B$ une $k$-algèbre, $S = \Spec(B)$ et $\omega$ un foncteur fibre de $\mathcal{T}$ sur $S$.

On fait $\mathcal{A} = \mathcal{T}$, $B_1 = B_2 = B$, $\omega_1 = \omega_2 = \omega$, (2.13)(i) permet d'appliquer 6.2. D'après 6.7, $\underline{\Aut}^k_k(\omega)$ est le spectre de $L(\omega) = L_k(\omega, \omega)$. L'action de $\underline{\Aut}^k_k(\omega)$ sur les $\omega(X)$ est définie par les morphismes
\[
\omega(X) \to \omega(X) \otimes L(\omega, \omega)
\]
qui définissent $L$ (4.7.3). Par 6.7 encore, la loi de composition du groupoïde $\underline{\Aut}^k_k(\omega)$ est définie par la comultiplication de $L(\omega)$ et 6.2 équivaut à 1.12(ii).

Passons à la preuve de 1.12(i) pour $S$ affine. Un point crucial sera le cas particulier du lemme suivant obtenu en faisant $I = \{1, 2\}$ et $\mathcal{T}_1 = \mathcal{T}_2 = \mathcal{T}$.

\begin{lemma}[6.9]
Soit $(\mathcal{T}_i)_{i \in I}$ une famille finie de catégories tannakiennes sur $k$. Le produit tensoriel sur $k$ $\bigotimes_k \mathcal{T}_i$ des $\mathcal{T}_i$, muni du produit tensoriel (5.16), est une catégorie tensorielle sur $k$.
\end{lemma}

Pour $k$ parfait, le résultat a déjà été prouvé (5.17). Le cas général sera traité en 6.16--6.20.

Si $\omega_i$ est un foncteur fibre de $\mathcal{T}_i$ sur $B_i \neq 0$, le foncteur exact de $\prod \mathcal{T}_i$ dans les $\otimes B_i$-modules: $(X_i) \mapsto \bigotimes \omega_i(X_i)$ se factorise par un foncteur $k$-linéaire exact à droite de $\bigotimes_k \mathcal{T}_i$ dans les $\otimes B_i$-modules. Étant exact à droite, c'est un foncteur fibre (2.10.1). Il résulte donc de 6.9 qu'un produit tensoriel de catégorie tannakiennes est encore une catégorie tannakienne.

\subsection{}

Soit $\mathcal{T}$ une catégorie tannakienne sur $k$. La catégorie abélienne $\mathcal{T}$ étant noethérienne (2.13(i)), tout Ind-objet de $\mathcal{T}$ est réunion croissante de ses sous-objets dans $\mathcal{T}$. Si $\omega$ est un foncteur fibre de $\mathcal{T}$ sur un $k$-schéma $S$, on étend $\omega$ en un $\otimes$-foncteur encore noté $\omega$ de $\Ind\mathcal{T}$ dans les modules quasi-cohérents sur $S$ par $\omega(\text{``}\varinjlim\text{''} X_\alpha) = \varinjlim \omega(X_\alpha)$. Pour $X$ dans $\Ind\mathcal{T}$, $\omega(X)$ est plat sur $S$ en tant que limite inductive de modules localement libres.

\begin{lemma}[6.11]
Avec les hypothèses et notations de 6.10, si le Ind-objet $X$ est non nul, $\omega(X)$ est fidèlement plat sur $S$.
\end{lemma}

\begin{proof}
Si $X \neq 0$, il existe $Y \hookrightarrow X$ avec $Y$ dans $\mathcal{T}$ et non nul. D'après 2.10.1, $\omega(Y)$ est localement libre partout de rang non nul, donc fidèlement plat sur $S$. On a $\omega(Y) \hookrightarrow \omega(X)$ avec un quotient $\omega(X/Y)$ plat. La fidèle platitude de $\omega(X)$ en résulte.
\end{proof}

\subsection{}

Soient $C$ une petite catégorie, $\mathcal{T}$ une catégorie tensorielle sur $k$ et $T_1, T_2$ deux foncteurs de $C$ dans $\mathcal{T}$. À l'imitation de 4.7, on se propose de définir un Ind-objet $\Lambda_\mathcal{T}(T_1, T_2)$ de $\mathcal{T}$, muni, pour tout $X$ dans $C$, de
\begin{equation}\label{eq:6.12.1}
T_1(X)^\vee \otimes T_2(X) \to \Lambda_\mathcal{T}(T_1, T_2)
\end{equation}
tel que pour tout $f : X \to Y$ dans $C$ le diagramme
\begin{equation}\label{eq:6.12.2}
\begin{tikzcd}
T_1(Y)^\vee \otimes T_2(Y) \arrow[r] \arrow[d, "T_1(f)^\vee \otimes \Id"'] & \Lambda_\mathcal{T}(T_1, T_2) \\
T_1(X)^\vee \otimes T_2(X) \arrow[ur]
\end{tikzcd}
\end{equation}
soit commutatif, et universel pour ces propriétés. Pour $C$ finie, $\Lambda$ est simplement un quotient convenable de la somme des $T_1(X)^\vee \otimes T_2(X)$ ($X$ dans $C$), et est dans $\mathcal{T}$. Dans le cas général, on écrit $C$ comme limite inductive filtrante de ses sous-catégories finies $C_\alpha$, et on pose
\[
\Lambda_\mathcal{T}(T_1, T_2) := \text{``}\varinjlim\text{''} \Lambda_\mathcal{T}(T_1|C_\alpha, T_2|C_\alpha).
\]

Soient $\mathcal{T}_1$ et $\mathcal{T}_2$ deux catégories tensorielles sur $k$ telles que le produit tensoriel sur $k$ $\mathcal{T} := \mathcal{T}_1 \otimes \mathcal{T}_2$ existe et soit une catégorie tensorielle sur $k$.

Soit $\inj_j$ l'injection de $\mathcal{T}_j$ dans $\mathcal{T}$ : $\inj_1 : X \mapsto X \otimes 1$, $\inj_2 : X \mapsto 1 \otimes X$. Pour $T_j$ un foncteur de $C$ dans $\mathcal{T}_j$, on pose
\[
\Lambda_k(T_1, T_2) := \Lambda_\mathcal{T}(\inj_1 \circ T_1, \inj_2 \circ T_2).
\]

Soient $B_1$ et $B_2$ des $k$-algèbres commutatives et $\omega_j$ un foncteur fibre de $\mathcal{T}_j$ sur $B_j$. Le foncteur $(X, Y) \mapsto \omega_1(X) \otimes_k \omega_2(Y)$ est un foncteur exact en chaque variable de $\mathcal{T}_1 \times \mathcal{T}_2$ dans les $B_1 \otimes_k B_2$-modules. Par définition du produit tensoriel $\mathcal{T}_1 \otimes \mathcal{T}_2$, il se factorise uniquement par un foncteur exact à droite
\[
\omega : \mathcal{T}_1 \otimes \mathcal{T}_2 \to (B_1 \otimes_k B_2\text{-modules}).
\]
Ce foncteur est un $\otimes$-foncteur. Étant exact à droite, il est un foncteur fibre (2.10.1). On a
\begin{equation}\label{eq:6.12.3}
\omega \Lambda_k(T_1, T_2) = L_k(\omega_1 T_1, \omega_2 T_2).
\end{equation}

Faisons $C = \mathcal{T}_1 = \mathcal{T}_2 = \mathcal{T}$ et prenons pour $T_i$ le foncteur identique de $\mathcal{T}$ dans $\mathcal{T}$. On obtient un ind-objet $\Lambda := \Lambda_k(\Id_\mathcal{T}, \Id_\mathcal{T})$ de $\mathcal{T} \otimes \mathcal{T}$. D'après (6.12.3), quels que soient les foncteurs fibres $\omega_1$ et $\omega_2$ de $\mathcal{T}$ sur $B_1$ et $B_2$, on a
\begin{equation}\label{eq:6.12.4}
\omega \Lambda = L_k(\omega_1, \omega_2).
\end{equation}

\begin{lemma}[6.13]
$\Lambda := \Lambda_k(\Id_\mathcal{T}, \Id_\mathcal{T})$ est non nul.
\end{lemma}

\begin{proof}
Soit $T$ le foncteur exact à droite de $\mathcal{T} \otimes \mathcal{T}$ dans $\mathcal{T}$ tel que $T(X \otimes Y) = X \otimes Y$. On a
\[
T\Lambda = \Lambda_\mathcal{T}(\Id_\mathcal{T}, \Id_\mathcal{T}).
\]
Les morphismes d'évaluation
\[
X^\vee \otimes X \to \mathbf{1}
\]
rendent commutatifs les diagrammes
\[
\begin{tikzcd}
X^\vee \otimes Y \arrow[r] \arrow[d] & X^\vee \otimes X \arrow[d] \\
Y^\vee \otimes Y \arrow[r] & \mathbf{1}
\end{tikzcd}
\]
pour $f : Y \to X$ et définissent $e : \Lambda_\mathcal{T}(\Id_\mathcal{T}, \Id_\mathcal{T}) \to \mathbf{1}$. Pour $X = \mathbf{1}$, (6.12.1) est un morphisme de $\mathbf{1} = \mathbf{1}^\vee \otimes \mathbf{1}$ dans $\Lambda_\mathcal{T}(\Id_\mathcal{T}, \Id_\mathcal{T})$. Le composé
\[
\mathbf{1} = \mathbf{1}^\vee \otimes \mathbf{1} \to \Lambda_\mathcal{T}(\Id_\mathcal{T}, \Id_\mathcal{T}) \to \mathbf{1}
\]
est l'identité. On a donc $\Lambda_\mathcal{T}(\Id_\mathcal{T}, \Id_\mathcal{T}) \neq 0$. A fortiori, $\Lambda \neq 0$.
\end{proof}

\subsection{Preuve de 1.12(i) pour $S$ affine}

Posons $S = \Spec(B)$. Avec les notations de 6.7, il s'agit de montrer que si $\omega$ un foncteur fibre de $\mathcal{T}$ sur $S$, $L_k(\omega, \omega)$ est fidèlement plat sur $B \otimes_k B$. Par (6.12.4) pour $\omega_1 = \omega_2 = \omega$, c'est une conséquence de 6.11 et 6.13.

La même preuve donne, directement plutôt que via la réduction 1.13, que pour $\omega_1$ et $\omega_2$ des foncteurs fibres sur $S_1$ et $S_2$, $\underline{\Isom}_k(\omega_1, \omega_2)$ est fidèlement plat sur $S_1 \times_k S_2$.

\subsection{Preuve de 1.12(iii) pour $S$ affine}

Soit $G$ un groupoïde agissant transitivement sur $S \neq 0$ et affine sur $S \times S$. Soit $\omega$ le foncteur fibre ``oubli de l'action de $G$'' de $\Rep(S : G)$ sur $S$. Il s'agit de prouver que $G \xrightarrow{\sim} \underline{\Aut}_k(\omega)$. Les deux membres étant des groupoïdes transitifs, il suffit de le prouver après passage à une fibre. Pour $S_1 \to S$, $S_1 \neq 0$, $G_1$ le groupoïde induit et $\omega_1$ le foncteur fibre oubli de $\Rep(S_1 : G_1)$ sur $S_1$, on a $\Rep(S : G) \xrightarrow{\sim} \Rep(S_1 : G_1)$ (3.5.1) et le groupoïde induit par $\underline{\Aut}_k(\omega)$ sur $S_1$ est donc $\underline{\Aut}_k(\omega_1)$. Prenant pour $S_1$ un point de $S$, ceci ramène à supposer que $S_1$ est le spectre d'un corps. Il ne reste qu'à appliquer 4.13.

Compte tenu de 1.13(b), ceci termine la preuve de 1.12.

\subsection{Preuve de 6.9}

Soit $\mathcal{T}$ une catégorie tensorielle sur $k$. Pour $X$ dans $\mathcal{T}$, la catégorie tensorielle $\langle X \rangle^\otimes$ engendrée par $X$ est la sous-catégorie pleine de $\mathcal{T}$ d'objets les sous-quotients des sommes de $X^{\otimes n} \otimes X^{\vee \otimes m}$. On dit que $\mathcal{T}$ est de $\otimes$-génération finie s'il existe $X$ dans $\mathcal{T}$ tel que $\mathcal{T} = \langle X \rangle^\otimes$. Il est clair que $\mathcal{T}$ est réunion filtrante de ses sous-catégories $\langle X \rangle^\otimes$ et, par passage à la limite inductive (5.1), il suffit de prouver 6.9 lorsque les $\mathcal{T}_i$ sont de $\otimes$-génération finie. L'assertion résulte de 5.21 et du lemme suivant.

\begin{lemma}[6.17]
Si $\mathcal{T}$ est une catégorie tannakienne sur $k$ de $\otimes$-génération finie, il existe un foncteur fibre $\omega$ de $\mathcal{T}$ sur $S$ spectre d'une extension finie de $k$ tel que $\underline{\Aut}_k(\omega)$ soit fidèlement plat sur $S \times S$ (d'où, par 1.12(ii) prouvé en 6.3, $\mathcal{T} \xrightarrow{\sim} \Rep(S : \underline{\Aut}_k(\omega))$).
\end{lemma}

Soient $\mathcal{T}$ comme en 6.17 et $\omega$ un foncteur fibre sur $S = \Spec(B) \neq 0$.

\begin{lemma}[6.18]
Le groupoïde $G := \underline{\Aut}_k(\omega)$ agissant sur $S$ est une limite projective de groupoïdes $G_n$ ($n > 0$) affines et de présentation finie sur $S \times S$, avec $G_m$ sous-schéma fermé de $G_n$ pour $m \geq n$.
\end{lemma}

\begin{proof}
Pour $V$ un faisceau localement libre de rang fini sur $S$, soit $\underline{\End}_k(V)$ (resp.~$\underline{\Aut}_k(V)$) le foncteur qui à $T$ sur $S \times S : (p_1, p_2) : T \to S \times S$ attache l'ensemble des morphismes (resp.~isomorphismes) $p_1^*V \to p_2^*V$. Le foncteur est représentable par un schéma affine sur $S \times S$ encore noté $\underline{\End}_k(V)$ (resp.~$\underline{\Aut}_k(V)$). Si $V$ est libre, le choix d'une base $e_1, \ldots, e_n$ de $V$ identifie $\underline{\Aut}_k(V)$ au fibré trivial sur $S \times S$ de fibre $\GL_n$. La composition des isomorphismes fait de $\underline{\Aut}_k(V)$ un groupoïde agissant sur $S$. Pour $V$ libre muni d'une base, c'est le produit du groupoïde $S \times S$ agissant sur $S$ et du groupe $\GL_n$.

Pour $Y$ dans $\mathcal{T}$, soit $\underline{\Aut}_k(\omega|\langle Y \rangle)$ le foncteur qui à $T$ sur $S \times S : (p_1, p_2) : T \to S \times S$ attache l'ensemble des isomorphismes de foncteurs $p_1^*(\omega|\langle Y \rangle) \xrightarrow{\sim} p_2^*(\omega|\langle Y \rangle)$. L'application de $\underline{\Aut}_k(\omega|\langle Y \rangle)(T)$ dans $\underline{\Aut}(\omega(Y))(T)$ est injective. Son image est l'ensemble des isomorphismes $g : \pr_1^*\omega(Y) \xrightarrow{\sim} \pr_2^*\omega(Y)$ tels que pour tout $Z \subset Y^n$, $g_{\otimes n} : \pr_1^*\omega(Y^{\otimes n}) = \pr_1^*\omega(Y)^{\otimes n} \xrightarrow{\sim} \pr_2^*\omega(Y)^{\otimes n} = \pr_2^*\omega(Y^{\otimes n})$ respecte les sous-modules $\pr_1^*\omega(Z)$. En effet, pour tout sous-quotient $Q = Z_1/Z_2$ de $Y^n$, un tel $g$ induit un isomorphisme $g_Q$ entre les $\pr_i^*\omega(Q) = \pr_i^*\omega(Z_1)/\pr_i^*\omega(Z_2)$, et un argument de graphe montre que ces $g_Q$ commutent aux $\omega(f)$, pour $f$ un morphisme entre sous-quotients de $Y^n$. Cette description montre que $\pr_i^*\omega(Q) = \pr_i^*\omega(Z_1)/\pr_i^*\omega(Z_2)$. Pour $Q'$ un sous-quotient $Z'_1/Z'_2$ de $Y^m$, et $f : Q \to Q'$, le graphe $\Gamma(f) \subset Q \times Q'$ de $f$ a une image inverse $f^{-1}(\Gamma) \subset Y^{n+m}$ respectée par $g_{\otimes n+m}$. On a donc $\omega(f) \circ g_Q = g_{Q'} \circ \omega(f)$ et $g$ se prolonge en le morphisme de foncteurs $(g_Q)$.

Cette description montre que $\underline{\Aut}_k(\omega|\langle Y \rangle)$ est représenté par un sous-schéma fermé --- encore noté $\underline{\Aut}_k(\omega|\langle Y \rangle)$ --- de $\underline{\Aut}_k(\omega(Y))$.

Si $P$ est un générateur projectif de $\langle Y \rangle$, $\underline{\Aut}_k(\omega|\langle Y \rangle)(T)$ s'identifie encore à l'ensemble des isomorphismes $g : \pr_1^*\omega(P) \xrightarrow{\sim} \pr_2^*\omega(P)$ qui commutent à l'action de $\End(P)$. Ceci montre que $\underline{\Aut}_k(\omega|\langle Y \rangle)$ est le sous-schéma de $\underline{\End}_k(\omega(P))$ défini par un nombre fini d'équations. C'est aussi le sous-schéma de $\underline{\End}_k(\omega(Y^n))$ défini par un nombre fini d'équations: si $P$ est le sous-quotient $Z_1/Z_2$ de $Y^n$, il s'agit de respecter les $\omega(Z_i) \subset \omega(Y^n)$, la structure de $\End P$-module de $\omega(P) = \omega(Z_1)/\omega(Z_2)$ et $\omega(f)$ pour $f$ un épimorphisme $P^m \to Y$.

Soit $X$ un $\otimes$-générateur de $\mathcal{T}$. Quitte à remplacer $X$ par $X \oplus \mathbf{1} \oplus X^\vee$, on peut supposer, et on suppose, que $\mathbf{1}$ et $X^\vee$ sont dans $\langle X \rangle$. Les $\mathcal{T}(X, n) := \langle X^{\otimes n} \rangle$ forment alors une suite croissante de sous-catégories de $\mathcal{T}$ de réunion $\mathcal{T}$. Soit $G_n$ le foncteur qui à $(p_1, p_2) : T \to S \times S$ attache l'ensemble des isomorphismes $g$ de $\pr_1^*|\mathcal{T}(X, n)$ avec $\pr_2^*\omega|\mathcal{T}(X, n)$ qui respectent les isomorphismes $\omega(Y \otimes Z) \to \omega(Y) \otimes \omega(Z)$ pour $Y$ dans $\mathcal{T}(X, i)$ et $Z$ dans $\mathcal{T}(X, j)$ avec $i + j \leq n$. Il suffit pour cela que $g$ commute aux images inverses par les $\pr_i^*$ des isomorphismes
\[
\omega(\otimes^m X) \to \omega(X)^{\otimes m}
\]
pour $m \leq n$. On en déduit que $G_n$ est représenté par un sous-schéma fermé de $\underline{\Aut}_k(\omega(X^{\otimes n}))$ défini par un nombre fini d'équations.

Le groupoïde $G$ s'identifie à l'intersection des sous-groupoïdes $G_n$ de $\underline{\Aut}_k(\omega(X))$ et 6.18 en résulte.
\end{proof}

\begin{proposition}[6.19]
Sous les hypothèses de 6.18, avec $B$ un corps, $G$ est fidèlement plat de présentation finie sur $S \times S$.
\end{proposition}

\begin{proof}
Avec les notations de 6.18, soient $G_\Delta$ (resp.~$G_n^\Delta$) la restriction de $G$ (resp.~$G_n$) à la diagonale $S \hookrightarrow S \times S$ de $S \times S$. C'est un schéma en groupe sur $S$. Puisque $S$ est noethérien, $G_\Delta = G_n^\Delta$ pour $n \geq n_0$.

Écrivons $B$ comme limite inductive filtrante de $k$-sous-algèbres de type fini $B_\alpha$ et soit $S_\alpha := \Spec(B_\alpha)$. Fixons $n$. Puisque $G_n$ est de présentation finie sur $S \times S$, des arguments standard montrent que pour $\alpha$ assez grand, $G_n$ provient d'un groupoïde de présentation finie $G_{n,\alpha}$ agissant sur $S_\alpha$.

Puisque $B_\alpha \subset B$, $S_\alpha$ est intègre. Remplaçant $S_\alpha$ par un ouvert affine dense, on peut supposer que les deux projections $p_i : G_{n,\alpha} \to S_\alpha$ sont plates et que la restriction $G_{n,\alpha}^\Delta$ de $G_{n,\alpha}$ à la diagonale $S_\alpha \hookrightarrow S_\alpha \times S_\alpha$ est plate sur $S_\alpha$. Sous ces hypothèses, le quotient au sens des faisceaux fppf $S_\alpha/G_{n,\alpha}$ est représentable par un espace algébrique $Q_\alpha$: on dispose d'un morphisme fidèlement plat $S_\alpha \to Q_\alpha$, égalisant la double flèche $G_{n,\alpha} \rightrightarrows S_\alpha$ et le morphisme $G_{n,\alpha} \to S_\alpha \times S_\alpha$ est fidèlement plat (3.11).

Un espace algébrique admet toujours un ouvert non vide qui est un schéma affine. Remplaçant $S_\alpha$ par l'image inverse d'un tel ouvert, on obtient une situation où $Q_\alpha = S_\alpha/G_{n,\alpha}$ est un schéma affine. On a le diagramme
\[
\begin{tikzcd}
G_n \arrow[r] \arrow[d] & G_{n,\alpha} \arrow[d] \\
S \arrow[r] \arrow[dr] & S_\alpha \arrow[d] \\
& Q_\alpha
\end{tikzcd}
\]

Pour $f$ l'image inverse sur $S$ d'un élément de $\Gamma(Q_\alpha, \mathcal{O})$, les deux images inverses de $f$ sur $G$ coïncident. D'après (2.1.4), on a $f \in k$ et $Q_\alpha = \Spec(k)$.

On conclut que $G_{n,\alpha}$ est fidèlement plat sur $S_\alpha \times S_\alpha$, donc $G_n$ fidèlement plat sur $S \times S$.

Pour $m \geq n \geq n_0$, le morphisme de groupoïdes agissant transitivement sur $S$ : $G_m \to G_n$ induit un isomorphisme $G_m^\Delta \xrightarrow{\sim} G_n^\Delta$. C'est donc un isomorphisme (3.5.2) et $G = G_n$ pour $n \geq n_0$, en particulier $G$ est fidèlement plat de présentation finie sur $S \times S$.
\end{proof}

\begin{corollary}[6.20]
Si $\mathcal{T}$ admet un $\otimes$-générateur et admet un foncteur fibre $\omega$ sur $S = \Spec(B) \neq 0$, alors $\mathcal{T}$ admet un foncteur fibre sur une extension finie convenable $k'$ de $k$.
\end{corollary}

\begin{proof}
Soit $s \in S$, de corps résiduel $K$. Par extension des scalaires, $\omega$ fournit un foncteur fibre sur $K$. On peut donc supposer, et on suppose, que $B$ est un corps. Écrivons $B$ comme limite inductive de sous-algèbres $B_\alpha$ de type fini sur $k$. D'après 6.19, le groupoïde $G$ agissant sur $S$ des automorphismes de $\omega$ est de présentation finie. Il provient d'un groupoïde de présentation finie $G_\alpha$ sur un $S_\alpha$. Puisque $G$ est fidèlement plat sur $S \times S$, pour $\beta$ assez grand, le groupoïde $G_\beta$ agissant sur $S_\beta$ induit par $G_\alpha$ est fidèlement plat sur $S_\beta \times S_\beta$ (EGA). Le changement de base de $S_\beta$ à $S$ est alors une équivalence: $\Rep(S_\beta : G_\beta) \to \Rep(S : G)$ (3.5.1). On déduit donc de $\omega$ un foncteur fibre sur $S_\beta$. La fibre de $\omega$ en un point fermé de $S_\beta$ est le foncteur fibre voulu.
\end{proof}

Le lemme 6.17 résulte de 6.19 et 6.20. Par 6.16, ceci termine la preuve de 6.9, et celle du théorème principal 1.12.

\subsection{Preuve de 5.18}

Soient $S_i$, $G_i$, $S$ et $G$ comme en 5.18. Soit $G_i^S$ le groupoïde agissant sur $S$ induit par $G_i$. On a
\[
\Rep(S_i : G_i^S) \xrightarrow{\sim} \Rep(S : G_i^S)
\]
et $G$ est le produit (sur $S \times S$) des $G_i^S$. Il s'agit donc de prouver que le produit tensoriel sur $S$:
\begin{equation}\label{eq:6.21.1}
\bigotimes_i \Rep(S : G_i^S) \to \Rep(S : G)
\end{equation}
fait de la catégorie but le produit tensoriel des catégories sources.

Changeons de notations et soient $G_i$ des $k$-groupoïdes agissant transitivement sur $S$, de produit $G$.

Soit $\mathcal{T}$ le produit tensoriel des $\mathcal{T}_i := \Rep(S : G_i)$.

Le produit tensoriel sur $S$ des foncteurs fibres $\omega_i$ ``oubli de l'action de $G_i^S$'' est un foncteur exact en chaque variable $\prod \mathcal{T}_i \to$ (modules sur $S$). Il se factorise par un foncteur fibre $\omega$ de $\mathcal{T}$ sur $S$. On a (1.12(ii)) $\mathcal{T} \xrightarrow{\sim} \Rep(S : \underline{\Aut}_k(\omega))$ et le morphisme défini par transport de structures $G = \prod_S G_i = \prod_S \underline{\Aut}_k(\omega_i) \to \underline{\Aut}_k(\omega)$ donne lieu à un diagramme essentiellement commutatif
\[
\begin{tikzcd}
\mathcal{T} \arrow[r] \arrow[d] & \Rep(S : G) \\
\Rep(S : \underline{\Aut}_k(\omega)) \arrow[ur]
\end{tikzcd}
\]
Il reste à prouver que le morphisme
\begin{equation}\label{eq:6.21.2}
G : \prod_{S \times S} \underline{\Aut}_k(\omega_i) \to \underline{\Aut}_k(\omega)
\end{equation}
est un isomorphisme. On dispose de $\otimes$-foncteurs $\inj_i : \mathcal{T}_i \to \mathcal{T} : X \mapsto$ produit tensoriel de $X$ avec les $\mathbf{1} \in \Ob\mathcal{T}_j$ ($j \neq i$). Ils définissent des morphismes $\underline{\Aut}_k(\omega_i) \to \underline{\Aut}_k(\omega_i)$, de produit une rétraction de (6.21.2). Cette rétraction est injective, car tout objet de $\mathcal{T}$ est sous-quotient d'un $\otimes \inj_i(X_i) = \otimes X_i$. Le morphisme (6.21.2) est donc un isomorphisme.

\section{Caractérisation interne des catégories tannakiennes (caractéristique 0)}

Notre but dans ce paragraphe est de prouver le théorème suivant.

\begin{theorem}[7.1]\label{thm:7.1}
Soit $\mathcal{T}$ une catégorie tensorielle (1.2, 2.1) sur $k$ de caractéristique 0. Les conditions suivantes sont équivalentes:

\textup{(i)} $\mathcal{T}$ est tannakienne (2.8).

\textup{(ii)} Pour tout $X$ dans $\mathcal{T}$, $\dim X$ est un entier $\geq 0$.

\textup{(iii)} Pour tout $X$ dans $\mathcal{T}$, il existe un entier $n \geq 0$ tel que $\bigwedge^n X = 0$.
\end{theorem}

Rappelons la définition de $\dim X$. Pour tout morphisme $f : X \to Y$ soit $\delta(f) : \mathbf{1} \to Y \otimes X^\vee$ le composé $(\Id_{X^\vee} \otimes f) \circ \delta : \mathbf{1} \to X^\vee \otimes X \to X^\vee \otimes Y = Y \otimes X^\vee$. Pour $X = Y$, on pose $\Tr(f) := \ev \circ \delta(f) \in \End(\mathbf{1}) = k$, et $\dim X := \Tr(\Id_X) = \ev \circ \delta$.

La puissance extérieure $\bigwedge^n X$ est l'image de l'antisymétrisation $a = \sum (-1)^{\varepsilon(\sigma)} \sigma : \otimes^n X \to \otimes^n X$. En caractéristique 0, c'est aussi l'image du projecteur $a/n!$, et
\begin{equation}\label{eq:7.1.1}
\dim \bigwedge^n X = \Tr(a/n!) = \Tr(a)/n!.
\end{equation}

Montrons que cette dimension est donnée par la formule habituelle
\begin{equation}\label{eq:7.1.2}
\dim \bigwedge^n X = \binom{d}{n} := (\dim X)(\dim X - 1) \cdots (\dim X - n + 1)/n!.
\end{equation}

\begin{lemma}[7.2]
Soient $n$ un entier $\geq 1$, $X_i$ ($i \in \mathbb{Z}/(n)$) dans $\mathcal{T}$, et des morphismes $u_i : X_i \to X_{i+1}$. Ils définissent un endomorphisme $\otimes u_i$ du produit tensoriel des $X_i$. La trace de cet endomorphisme coïncide avec celle du composé des $u_i$:
\begin{equation}\label{eq:7.2.1}
\Tr(\otimes u_i) = \Tr(u_n \circ \cdots \circ u_1).
\end{equation}
\end{lemma}

\begin{proof}
Pour le composé de deux morphismes $X \xrightarrow{u} Y \xrightarrow{v} Z$, $\delta(vu)$ est le composé
\[
\mathbf{1} \xrightarrow{\delta(u) \otimes \Id_{X^\vee}} Z \otimes Y^\vee \otimes Y \otimes X^\vee \xrightarrow{\Id \otimes \ev_Y} Z \otimes X^\vee.
\]
Itérant, on obtient que $\delta(u_n \circ \cdots \circ u_1)$ est le composé
\[
\begin{tikzcd}
\mathbf{1} \arrow[r, "\delta(u_n) \otimes \cdots \otimes \delta(u_1)"'] & X_1 \otimes X_n^\vee \otimes X_n \otimes X_{n-1}^\vee \otimes \cdots \otimes X_2 \otimes X_1^\vee \arrow[d, "\ev(X_n) \otimes \cdots \otimes \ev(X_2)"'] \\
& X_1 \otimes X_1^\vee
\end{tikzcd}
\]
et la trace de $u_n \cdots u_1$ est le composé
\[
\delta(u_n \cdots u_1) = \ev \circ (\otimes \delta(u_i)).
\]
La trace de $\otimes u_i$ est donnée par la même formule et 7.2 en résulte.
\end{proof}

Prouvons (7.1.2). Si $\sigma \in \mathfrak{S}_n$ est une permutation cyclique, 7.2 appliqué à $X_i = X$, $u_i = \Id_X$ donne
\[
\Tr(\sigma, \otimes^n X) = \dim(X).
\]
Pour $\sigma$ une permutation ayant $k$ cycles, on en déduit
\[
\Tr(\sigma, \otimes^n X) = \dim(X)^k
\]
et (7.1.1) prouve l'existence d'un polynôme universel $P \in \mathbb{Q}[T]$ tel que
\[
\dim \bigwedge^n X = P(\dim X).
\]
Prenant pour $\mathcal{T}$ la catégorie des espaces vectoriels sur $k$, on trouve que pour tout entier $d \geq 0$,
\[
P(d) = \binom{d}{n}.
\]
On a donc $P(T) = \binom{T}{n}$ et (7.1.2) en résulte.

\begin{lemma}[7.3]
Si $\mathcal{T}$ vérifie 7.1\textup{(ii)} et que $X$ dans $\mathcal{T}$ est de dimension 0, alors $X = 0$.
\end{lemma}

\begin{proof}
Si $X \neq 0$, l'application identique de $X$ est non nulle et $\delta : \mathbf{1} \to X \otimes X^\vee$ est non nul, donc injectif (cf.~2.9). On a $\dim(X \otimes X^\vee) = \dim X \cdot \dim X^\vee = 0$ et $\dim \coker(\delta) = \dim(X \otimes X^\vee) - \dim \mathbf{1} = -1$: contradiction.
\end{proof}

\subsection{Preuve de \ref{thm:7.1}: (ii) $\Rightarrow$ (iii)}

Si $\dim X$ est un entier $n \geq 0$, (7.1.2) donne $\dim \bigwedge^{n+1} X = 0$, donc par 7.3 $\bigwedge^{n+1} X = 0$.

\subsection{Preuve de \ref{thm:7.1}: (iii) $\Rightarrow$ (ii)}

Si $\bigwedge^n X = 0$, on a $\dim \bigwedge^n X = 0$ et (7.1.2) montre que $\dim X$ est un entier entre 0 et $n-1$.

Il est clair que (i) $\Rightarrow$ (ii)(iii). Il reste à montrer que si $\mathcal{T}$ vérifie (ii) et (iii), alors $\mathcal{T}$ est tannakienne. La preuve sera donnée en 7.18.

\subsection{}

Dans les $n^\circ$ 7.5 à 7.12, nous allons montrer comment on peut imiter, dans une catégorie tensorielle quelconque $\mathcal{T}$, des rudiments d'algèbre commutative et de géométrie algébrique. Nous n'exigeons pas que $k$ soit de caractéristique 0.

Soit $\Ind\mathcal{T}$ la catégorie des Ind-objets de $\mathcal{T}$ (SGA4 I 8.2). Une définition est: un objet est un système inductif filtrant $X_i$, et les flèches sont définies par
\[
\Hom((X_i), (Y_j)) = \varinjlim_i \varprojlim_j \Hom(X_i, Y_j).
\]
Le système inductif filtrant $(X_i)$, vu comme objet de $\Ind\mathcal{T}$, est noté ``$\varinjlim$'' $X_i$.

Une définition équivalente: la catégorie des foncteurs $\mathcal{T} \to (\Ens)$, limites inductives filtrantes de foncteurs représentables. Au système inductif $(X_i)$, associer la limite inductive des foncteurs $h_{X_i}$. Parce que $\mathcal{T}$ est une catégorie abélienne, $\Ind(\mathcal{T})$ l'est aussi. Le produit tensoriel dans $\mathcal{T}$ en définit un dans $\Ind\mathcal{T}$, encore exact, associatif et commutatif.

Un anneau (à unité) $A$ de $\Ind\mathcal{T}$ est un objet $A$ de $\Ind\mathcal{T}$ muni d'un produit $A \otimes A \to A$ associatif et admettant une ``unité'' $\mathbf{1} \to A$ (un morphisme $u$ tels que les composés
\[
\begin{tikzcd}
A \arrow[r, "u \otimes \Id_A"'] & A \otimes A \arrow[r] & A
\end{tikzcd}
\quad\text{et}\quad
\begin{tikzcd}
A \arrow[r, "\Id_A \otimes u"'] & A \otimes A \arrow[r] & A
\end{tikzcd}
\]
soient l'identité).

On définit de façon évidente les $A$-modules à gauche et à droite, le produit tensoriel sur $A$, la commutativité de $A$. Par exemple: un $A$-module à gauche $M$ est un objet de $\Ind\mathcal{T}$ muni d'un produit $A \otimes M \to M$ rendant commutatif le diagramme
\[
\begin{tikzcd}
A \otimes A \otimes M \arrow[r] \arrow[d] & A \otimes M \arrow[d] \\
A \otimes M \arrow[r] & M
\end{tikzcd}
\]
et tel que le composé
\[
\begin{tikzcd}
M \arrow[r, "u \otimes \Id_M"'] & \mathbf{1} \otimes M \arrow[r] & A \otimes M \arrow[r] & M
\end{tikzcd}
\]
soit l'identité. On a $M \otimes_A N = \coker(M \otimes A \otimes N \rightrightarrows M \otimes N)$. Pour $A$ commutatif, la catégorie des $A$-modules, munie du produit tensoriel sur $A$, vérifie (2.1.1). L'objet unité est le $A$-module $A$, encore noté $\mathbf{1}_A$.

\begin{example}[7.6]
\textup{(i)} L'objet 0 admet une unique structure d'anneau. Tout module sur cet anneau est nul.

\textup{(ii)} L'objet unité $\mathbf{1}$, muni du morphisme structural $\mathbf{1} \otimes \mathbf{1} \to \mathbf{1}$ et de l'identité $\mathbf{1} \to \mathbf{1}$ est un anneau. Un objet de $\Ind(\mathcal{T})$ admet exactement une structure de module sur cet anneau.

L'anneau $\mathbf{1}$ (resp.~0) est initial (resp.~final) dans la catégorie des anneaux de $\Ind(\mathcal{T})$.
\end{example}

\subsection{}

Soient $f : A \to B$ un morphisme d'anneaux commutatifs de $\Ind(\mathcal{T})$. On dit que $B$ est plat (resp.~fidèlement plat) sur $A$ si le foncteur ($A$-modules) $\to$ ($B$-modules) : $M \mapsto B \otimes_A M$ est exact (resp.~exact et fidèle). Le formalisme de la descente fidèlement plate des modules (SGA 1 VIII 1) s'applique tel quel: si $B$ est fidèlement plat sur $A$, le foncteur $M \mapsto B \otimes_A M$ est une équivalence de la catégorie des $A$-modules avec celle de $B$-modules $N$ munis d'une donnée de descente (un isomorphisme $(B \otimes_A B) \otimes_B N \xrightarrow{\sim} N \otimes_B (B \otimes_A B)$ vérifiant la condition de cocycle usuelle). La preuve de SGA 1 VIII 1 s'applique, ou on peut se ramener au théorème de Barr-Beck (4.1) comme en 4.2.

Si $A \neq 0$, le morphisme unité $\mathbf{1} \to A$ est non nul, donc un monomorphisme (cf.~2.9). Le tensorisant avec un Ind-objet $M$, on obtient un monomorphisme $M \hookrightarrow A \otimes M$. Si $A \neq 0$, on a donc $A \otimes M \neq 0$. Le produit tensoriel dans $\Ind\mathcal{T}$ coïncide avec le produit tensoriel de modules sur l'anneau $\mathbf{1}$. Le foncteur $M \mapsto A \otimes_\mathbf{1} M = A \otimes M$ est donc exact et fidèle. Si $A \neq 0$, $A$ est donc fidèlement plat sur $\mathbf{1}$.

\subsection{}

Pour disposer d'un langage plus géométrique, on appelle \emph{catégorie des schémas affines en $\mathcal{T}$} l'opposée de celle des anneaux commutatifs à unité de $\Ind(\mathcal{T})$. On dira aussi: $\mathcal{T}$-schéma affine. On note $\Sp(A)$ le $\mathcal{T}$-schéma affine défini par $A$. Les produits fibrés existent: ils correspondent au produit tensoriel. Les sommes disjointes finies existent: elles correspondent aux produits.

Soit $S$ un $\mathcal{T}$-schéma affine. Un $S$-schéma affine est un $\mathcal{T}$-schéma affine $X$, muni de $a : X \to S$. Un $S$-schéma en groupes affines est un objet en groupes de la catégorie ($\mathcal{T}$-schémas affines/$S$) des $\mathcal{T}$-schémas affines sur $S$.

On appelle schéma vide (resp.~point noté $(\mathrm{Pt})$) le schéma initial (resp.~final) $\Sp(0)$ (resp.~$\Sp(\mathbf{1})$). On dit que $S = \Sp(A)$ est non vide si $A \neq 0$.

Puisque $(\mathrm{Pt})$ est final, il n'y a pas lieu de distinguer entre $(\mathrm{Pt})$-schémas affines et $\mathcal{T}$-schémas affines, ni entre $(\mathrm{Pt})$-schémas en groupes affines et objets en groupes de la catégorie des $\mathcal{T}$-schémas affines, i.e.~algèbres de Hopf commutatives à antipode de $\Ind(\mathcal{T})$.

Pour abréger, nous dirons simplement $S$\emph{-groupe} pour $S$-schéma en groupes affines et $\mathcal{T}$\emph{-groupe} pour $(\mathrm{Pt})$-schéma en groupes affines.

On dit que $T = \Sp(B)$ est fidèlement plat sur $S = \Sp(A)$ si $B$ est fidèlement plat sur $A$. On dit que des $S_i \to S$ ($i \in I$) couvrent $S$ s'il existe $T$ fidèlement plat sur $S$, une partie finie $J$ de $I$ et un diagramme commutatif
\[
\begin{tikzcd}
T \arrow[r] \arrow[d] & \coprod_{i \in J} S_i \arrow[dl] \\
S
\end{tikzcd}
\]
Ceci définit une topologie sur la catégorie des $\mathcal{T}$-schémas affines. L'adjectif ``localement'' référera à cette topologie.

\textit{Exemple.} Si le $\mathcal{T}$-schéma $S$ est non vide, il est fidèlement plat sur le point.

Comme en théorie des schémas, il est souvent commode de décrire un $S$-schéma affine $X$ en terme du foncteur représentable $T/S \mapsto \Hom_S(T, X)$ correspondant. Ce foncteur est un faisceau pour la topologie fidèlement plate. On appellera $\Hom_S(T, X)$ l'ensemble des $T$-points de $X$. Comme d'habitude, se donner une structure de $S$-groupe sur $X$ revient à se donner une structure fonctorielle de groupe sur l'ensemble de ses $T$-points, pour tout $T/S$. Cas particulier $S = (\mathrm{Pt})$: l'ensemble $X(T)$ des $T$-points de $X$ est $\Hom(T, X)$.

Si $S = \Sp(A)$, un $A$-module $M$ sera encore appelé un module sur $S$. Le produit tensoriel est le produit tensoriel sur $A$. Il vérifie (2.1.1). On écrit $\mathbf{1}_S$ pour $\mathbf{1}_A$ (le $A$-module unité $A$). Les catégories de modules forment une catégorie fibrée sur celle des $\mathcal{T}$-schémas affines. Pour $\Sp(B)$ au-dessus de $\Sp(A)$, l'image inverse de $M$ sur $\Sp(A)$ est $B \otimes_A M$ sur $\Sp(B)$. Le formalisme de la descente fidèlement plate pour les schémas affines s'applique.

Pour $M$ dans $\Ind(\mathcal{T})$, posons $F(M) := \Hom(\mathbf{1}, M)$. Pour $M$ un module sur $S = \Sp(A)$, on a
\[
F(M) = \Hom(\mathbf{1}, \mathbf{1}_A) \simeq \Hom_A(\mathbf{1}_A, M)
\]
et on appelle encore $F(M)$ les \emph{sections globales} de $M$ sur $S$. Prendre garde que le foncteur $F$ peut ne pas être exact.

\begin{example}[$V(M)$, 7.9]
Soit $M$ un module sur $S = \Sp(A)$. La puissance symétrique $\Sym^n_A(M)$ est le plus grand quotient de $\otimes^n_A M$ sur lequel le groupe symétrique $\mathfrak{S}_n$ agisse trivialement. L'\emph{algèbre symétrique} $\Sym^*_A(M)$ est la somme des $\Sym^n_A(M)$ ($n \geq 0$). Pour toute $A$-algèbre commutative $B$, on a
\[
\Hom_{A\text{-alg}}(\Sym^*_A(M), B) \simeq \Hom_A(M, B).
\]
Pour tout $T = \Sp(B)$ au-dessus de $S$, soit $M_T$ déduit de $M$ par changement de base. Le foncteur $T \mapsto \Hom(M_T, \mathbf{1}_T)$ est représentable, représenté par $V(M) := \Sp(\Sym^*_A(M))$: on a
\[
\Hom_B(B \otimes_A M, B) \simeq \Hom_A(M, B) \simeq \Hom_{A\text{-alg}}(\Sym^*_A(M), B).
\]
Le foncteur $T \mapsto \Hom(M_T, \mathbf{1}_T)$ est un foncteur en groupes et $V(M)$ est donc un $S$-groupe.
\end{example}

\begin{example}[Fibrés vectoriels, 7.10]
Si $M$ sur $S = \Sp(A)$ admet un dual $M^\vee$ (2.2 pour $M$ la catégorie des $A$-modules), tout facteur direct de $M$ en admet un également et, pour $T$ sur $S$, $(M^\vee)_T$ est encore un dual de $M_T$. En particulier, si $M$ est facteur direct d'un $A \otimes X$ avec $X$ dans $\mathcal{T}$, $M$ admet un dual.

Si $M$ et $N$ admettent un dual, on dispose (2.4) de
\[
u \mapsto {}^t u : \Hom_A(M, N) \xrightarrow{\sim} \Hom_A(N^\vee, M^\vee).
\]
Pour $N = \mathbf{1}_A$, on trouve en appliquant 7.9 que le foncteur qui à $T$ sur $S$ attache l'ensemble des sections globales sur $T$ de $M_T$ est représentable, représenté par $V(M^\vee)$:
\[
F(M_T) = \Hom_T(\mathbf{1}_T, M_T) = \Hom_T(M_T^\vee, \mathbf{1}_T).
\]
On appellera encore le $\mathcal{T}$-schéma $V(M^\vee)$ le \emph{fibré vectoriel} $M$.

Pour $S = (\mathrm{Pt})$, et pour $M$ dans $\mathcal{T}$, on appellera de même $\Sp(\Sym^*(M^\vee))$ le $\mathcal{T}$-schéma vectoriel $M$. Cette terminologie est parallèle à celle qui identifie un espace vectoriel $V$ de dimension finie sur $k$ et le schéma $\Spec \Sym^*(V^\vee)$ ayant même $k$-points.
\end{example}

\begin{example}[$k$-schémas et $\mathcal{T}$-schémas, $\mathbb{G}_a$, $\mathbb{G}_m$, 7.11]
Si $\eta : \mathcal{T}_1 \to \mathcal{T}_2$ est un $\otimes$-foncteur entre catégories tensorielles sur $k$, il s'étend en un $\otimes$-foncteur encore noté $\eta : \Ind\mathcal{T}_1 \to \Ind\mathcal{T}_2$. Ce dernier transforme anneaux de $\Ind\mathcal{T}_1$ en anneaux de $\Ind\mathcal{T}_2$, et $\mathcal{T}_1$-schémas affines en $\mathcal{T}_2$-schémas affines.

Prenons pour $\mathcal{T}_2$ une catégorie tensorielle $\mathcal{T}$, pour $\mathcal{T}_1$ la catégorie des multiples de $\mathbf{1}$, identifiée à celle des $k$-vectoriels de dimension finie (2.9) et pour $\eta$ l'inclusion de $\mathcal{T}_1$ dans $\mathcal{T}$. On trouve qu'un schéma affine sur $k$ définit un $\mathcal{T}$-schéma affine. La droite affine sur $k$ définit ainsi le schéma vectoriel $\mathbf{1}$, qu'on notera $\mathbb{A}^1$; son image inverse sur $S = \Sp(A)$ est le fibré vectoriel $\mathbf{1}_S$, qu'on notera $\mathbb{A}^1_S$. On a $\mathbb{A}^1 = \Sp(\Sym^*_\mathbf{1}(\mathbf{1}_A))$ et $\Sym^*_A(\mathbf{1}_A) = \bigoplus_{n \geq 0} A = A[X]$. L'ensemble $\Hom_S(S, \mathbb{A}^1_S)$ des sections de $\mathbb{A}^1_S$ est $F(A)$. On notera $1$ la section unité.

De même, le $k$-schéma $\mathbb{G}_m$ définit le $\mathcal{T}$-schéma qui représente le foncteur $S = \Sp(A) \mapsto F(A)^\times$.
\end{example}

\begin{example}[Espace homogène sous un vectoriel, 7.12]
Soit dans $\mathcal{T}$ une suite exacte
\[
0 \to M \to E \xrightarrow{u} N \to 0.
\]
Soit $v : \mathbf{1} \to E^\vee$ le transposé de $u$. Le morphisme $v$ définit un morphisme de $\mathcal{T}$-schémas vectoriels $E \to \mathbb{A}^1$. Soit $P$ le $\mathcal{T}$-schéma image inverse du point 1. C'est le produit fibré
\[
\begin{tikzcd}
P \arrow[r] \arrow[d] & E \arrow[d, "v"] \\
(\mathrm{Pt}) \arrow[r, "1"'] & \mathbb{A}^1
\end{tikzcd}
\]
i.e.~le $\Sp$ du produit tensoriel $\Sym^*(E^\vee) \otimes_{\Sym^*(\mathbf{1})} \mathbf{1}$. Les flèches sont $\Sym^*(\mathbf{1}) \to \mathbf{1}$: en degré 0, l'identité et $\Sym^1(\mathbf{1}) \to \Sym^1 E^\vee$ : en degré 1, $v$.

Le produit tensoriel est la limite inductive des $\Sym^n(E^\vee)$, avec pour morphismes de transition la multiplication par $v$:
\[
P = \Sp(\varinjlim \Sym^n(E^\vee)).
\]

Supposons $k$ de caractéristique 0. La puissance symétrique $\Sym^n X$ de $X$ dans $\mathcal{T}$ est alors le facteur direct de $\otimes^n X$ défini par la projection $s/n!$, avec $s = \sum \sigma$ ($\sigma \in \mathfrak{S}_n$). Si $X_1 \subset X$, le produit tensoriel des filtrations $0 \subset X_1 \subset X$ est une filtration sur $\otimes^n X$. Parce que le produit tensoriel est exact, le gradué associé est $\otimes^n (X_1 \oplus X/X_1)$. Passant au facteur direct défini par $s/n!$, on obtient une filtration de $\Sym^n(X)$, de quotients successifs les $\Sym^p(X_1) \otimes \Sym^{n-p}(X/X_1)$.

Dans le cas de la suite exacte
\[
\begin{tikzcd}
0 \arrow[r] & \mathbf{1} \arrow[r] & E^\vee \arrow[r] & M^\vee \arrow[r] & 0
\end{tikzcd}
\]
on obtient une filtration croissante $F$ de $\Sym^n(E^\vee)$ de quotients successifs les $\Sym^p(M^\vee)$. On a $F_p = \Sym^p(E^\vee)$ ($p \leq n$):
\[
0 \subset \Sym^0(E^\vee) \subset E^\vee = \Sym^1(E^\vee) \subset \Sym^2(E^\vee) \subset \cdots \subset \Sym^n(E^\vee).
\]
\end{example}

\subsection{}

Soit $\mathcal{P}$ une propriété de modules, stable par changements de bases. En accord avec 7.8, on dira qu'un module $M$ sur $S$ vérifie $\mathcal{P}$ localement s'il existe des $S_i \to S$ couvrant $M$ tel que les $M_i$ sur $S_i$ déduis de $M$ par changement de base vérifient $\mathcal{P}$. La propriété $\mathcal{P}$ est dite locale si elle est vraie dès que localement vraie.

\textit{Exemple:} la propriété d'admettre un dual est locale.

\begin{lemma}[7.14]
Supposons $k$ de caractéristique 0 et soit $0 \to M \to E \to N \to 0$ une suite exacte dans $\mathcal{T}$. Localement, cette suite est scindable.
\end{lemma}

\begin{proof}
Supposons tout d'abord que $N = \mathbf{1}$. Soit $P$ comme en 7.12. Parce que $k$ est de caractéristique 0, les calculs 7.12 donnent $P \neq 0$. Le $\mathcal{T}$-schéma $P$ est donc fidèlement plat sur le point. Après changement de base à $P$, il existe par définition une section globale de $E_P$ de projection sur $N_P = \mathbf{1}_P$ la section unité. Une telle section globale n'est pas autre chose qu'un scindage de la suite exacte $0 \to M_P \to E_P \to \mathbf{1}_P \to 0$ de modules sur $P$.

Passons au cas général. Prenons l'image inverse de l'extension $0 \to \underline{\Hom}(N, M) \to \underline{\Hom}(N, E) \to \underline{\Hom}(N, N) \to 0$ par $\delta : \mathbf{1} \to \underline{\Hom}(N, N)$. Il revient au même de scinder $(M, E, N)$ ou cette nouvelle extension; ceci nous ramène au cas déjà traité.
\end{proof}

Si un module $M$ sur $S = \Sp(A)$ admet un dual, sa dimension $\dim_A(M) \in F(A)$ est définie comme le composé $\ev \circ \delta : \mathbf{1}_A \to M^\vee \otimes M \to \mathbf{1}_A$.

\begin{lemma}[7.15]
Soit $M$ un module sur $S = \Sp(A)$. On suppose que $k$ est de caractéristique 0, que $A \neq 0$, que $M$ est facteur direct d'un $A$-module $A \otimes X$ ($X$ dans $\mathcal{T}$) et que $\dim_A M$ est un entier $d > 0$. Alors, $M$ admet localement une décomposition $M = \mathbf{1}_S \oplus N$.
\end{lemma}

\begin{proof}
La donnée d'une décomposition $M = \mathbf{1}_S \oplus N$ équivaut à celle de morphismes $\mathbf{1}_S \xrightarrow{v} M \xrightarrow{u} \mathbf{1}_S$ avec $uv = 1$. Soit $F$ le foncteur qui à $T$ sur $S$ attache l'ensemble des paires $v, u : \mathbf{1}_T \xrightarrow{v} M_T \xrightarrow{u} \mathbf{1}_T$ avec $vu = 1$.

Nous allons montrer que $F$ est représentable, représenté par un $\mathcal{T}$-schéma affine $\Sp(B)$ sur $\Sp(A)$, et que $B$ est fidèlement plat sur $A$.

Le foncteur des paires $v, u : \mathbf{1}_T \xrightarrow{v} M_T \xrightarrow{u} \mathbf{1}_T$ est représenté par le fibré vectoriel $M \times M^\vee$ sur $S$. La composition $(v, u) \mapsto vu$ est un morphisme de foncteur, définissant un morphisme de $\mathcal{T}$-schémas $c : M \times M^\vee \to \mathbb{A}^1_S$. Le foncteur $F$ est représenté par le produit fibré
\[
\begin{tikzcd}
\Sp(B) \arrow[r] \arrow[d] & M \times M^\vee \arrow[d, "c"] \\
(\mathrm{Pt}) \arrow[r, "1"'] & \mathbb{A}^1_S
\end{tikzcd}
\]

Calculons $B$. On a $\mathbb{A}^1_S = \Sp \Sym^*_A(\mathbf{1}_A)$, $\Sym^*_A(\mathbf{1}_A) = A[T] = \bigoplus_{n \geq 0} A$ et $(\mathrm{Pt}) \to \mathbb{A}^1_S$ est $T \mapsto 1 : A[T] \to A$ i.e.~l'identité en chaque degré:
\[
\Sym^n_A(\mathbf{1}_A) \to A.
\]
On a $M \times M^\vee = \Sp(\Sym^*_A(M^\vee) \otimes \Sym^*_A(M))$ et la composante de degré un, $A$, de $\Sym^*_A(\mathbf{1}_A)$ s'envoie dans $M^\vee \otimes M$; la flèche est $\delta$. Le produit tensoriel $B$ est le plus grand anneau quotient de $\Sym^*_A(M^\vee) \otimes \Sym^*_A(M)$ dans lequel les flèches unité et $\delta : A \to M^\vee \otimes M$ sont égalisées. C'est
\begin{equation}\label{eq:7.15.1}
B = \varinjlim_{p \geq 0} \Sym^n_A(M^\vee) \otimes_A \Sym^n_A(M)
\end{equation}
les flèches de transition dans la limite inductive étant la multiplication par $\delta$.

La puissance symétrique $\Sym^n_A M$ est plate sur $A$: si $M$ est facteur direct de $A \otimes X$, c'est un facteur direct du $A$-module plat $A \otimes \Sym^n X = \Sym^n(A \otimes X)$. La formule (7.15.1) montre donc que, comme $A$-module, $B$ est somme de limites inductives filtrantes de modules plats sur $A$.

Pour conclure, nous allons montrer que le morphisme $A \to B$ admet une rétraction $A$-linéaire. Il suffit pour cela de construire un système de morphismes
\[
r_n : \Sym^n_A(M^\vee) \otimes_A \Sym^n_A(M) \to A
\]
avec $r_0 = \Id$, rendant commutatifs les diagrammes
\begin{equation}\label{eq:7.15.2}
\begin{tikzcd}
\Sym^n_A(M^\vee) \otimes \Sym^n_A(M) \arrow[r, "\delta"] \arrow[dr, "r_n"'] & \Sym^{n+1}_A(M^\vee) \otimes \Sym^{n+1}_A(M) \arrow[d, "r_{n+1}"] \\
& A
\end{tikzcd}
\end{equation}
Parce qu'on est en caractéristique 0, la puissance symétrique $\Sym^n_A$ est le facteur direct de $\otimes^n$ défini par le projecteur $s/n!$ (cf.~7.12) et la dualité entre $\otimes^n M^\vee$ et $\otimes^n M$ induit une dualité entre $\Sym^n_A(M^\vee)$ et $\Sym^n_A(M)$. On prend
\begin{equation}\label{eq:7.15.3}
r_n := \ev / \dim_A \Sym^n M
\end{equation}
où $\dim_A \Sym^n M = d(d+1) \cdots (d+n-1)/n!$ (cf.~7.1).
\end{proof}

\begin{lemma}[7.16]
Pour $r_n$ donné par (7.15.3), les diagrammes (7.15.2) commutent.
\end{lemma}

\begin{proof}
La multiplication par $\delta$ de (7.15.2) se déduit de l'opération analogue dans $\otimes$ par symétrisation:
\[
\begin{tikzcd}
\Sym^n_A(M^\vee) \otimes \Sym^n_A(M) \arrow[r] \arrow[d, "s/(n+1)! \otimes s/(n+1)!"'] & \Sym^{n+1}_A(M^\vee) \otimes \Sym^{n+1}_A(M) \\
\otimes^{n+1} M^\vee \otimes \otimes^{n+1} M \arrow[r] & \otimes^{n+1} M^\vee \otimes \otimes^{n+1} M \arrow[u]
\end{tikzcd}
\]
On a $\ev \circ (s/(n+1)! \otimes s/(n+1)!) = \ev \circ (\Id \otimes s/(n+1)!)$. Développons en $\sigma \in \mathfrak{S}_{n+1}$. Le terme relatif à $\sigma$ ne dépend que de si $\sigma(n+1) = n+1$ ou non. On obtient
\[
\ev \circ (\delta \cdot) = \frac{1}{n+1} \cdot n \cdot (\ev, \delta) \cdot \ev \quad \text{(termes où $\sigma(n+1) = n+1$)}
\]
\[
+ \frac{n+1}{n+1} \cdot \ev \quad \text{(autres termes)}
\]
\[
= \frac{d+n}{n+1} \cdot \ev
\]
d'où 7.16.
\end{proof}

\begin{remark}[7.16.1]
Dans la catégorie tensorielle sur $k$ des super espaces vectoriels de dimension finie (1.4) un espace vectoriel $V$ purement impair n'admet pas localement $\mathbf{1}$ pour facteur direct. Le lecteur vérifiera que dans ce cas le foncteur $F$ de la preuve de 7.15 est le foncteur vide.
\end{remark}

\begin{proposition}[7.17]
Sous les hypothèses \textup{(ii)(iii)} de \ref{thm:7.1}, tout $X$ dans $\mathcal{T}$ est localement isomorphe à $\mathbf{1}^{\dim X}$.
\end{proposition}

\begin{proof}
Soit $d = \dim X$. Appliquant 7.15 de façon itérée, on obtient $A \neq 0$ et un isomorphisme de $A$-modules entre $A \otimes X$ et un $A$-module $\mathbf{1}_A^d \oplus N$. On a pour les $A$-modules
\[
\bigwedge^n_A(\mathbf{1}_A^d \oplus N) = \bigoplus_{p+q=n} \bigwedge^p_A(\mathbf{1}_A^d) \otimes \bigwedge^q_A(N).
\]
Pour $n = d+1$, on a $\bigwedge^{d+1}_A \mathbf{1}_A^d = 0$ et $N$ est donc facteur direct de
\[
\bigwedge^{d+1}_A(\mathbf{1}_A^d \oplus N) \simeq \bigwedge^{d+1}_A(A \otimes X) = A \otimes \bigwedge^{d+1} X = 0.
\]
On a donc $N = 0$, ce qui prouve 7.17.
\end{proof}

\subsection{Preuve de \ref{thm:7.1} (fin)}

Appliquant 7.14 et 7.17 transfiniment, on obtient un anneau commutatif $A \neq 0$ de $\Ind\mathcal{T}$ tel que
\begin{enumerate}
\item[(7.19.1)] Pour tout $X$ dans $\mathcal{T}$, les $A$-modules $A \otimes X$ et $\mathbf{1}_A^{\dim X}$ sont isomorphes;
\item[(7.19.2)] Pour toute suite exacte courte dans $\mathcal{T}$, la suite exacte courte de $A$-modules qui s'en déduit par extension des scalaires est scindable.
\end{enumerate}

Le foncteur $X \mapsto F(A \otimes X)$ est alors un foncteur fibre sur $F(A)$.

\section{Le groupe fondamental d'une catégorie tensorielle}

\subsection{}

Dans ce paragraphe, nos résultats essentiels supposent que $\mathcal{T}$ est une catégorie tensorielle sur $k$ vérifiant la condition de finitude (2.12.1): objets de longueur finie et $\Hom$ de dimension finie, et que $\mathcal{T} \otimes \mathcal{T}$ est encore une catégorie tensorielle sur $k$. Cette dernière condition est vérifiée pour $k$ parfait (5.17) ou pour $\mathcal{T}$ tannakienne (6.9). J'ignore si elle est toujours vérifiée.

L'idée essentielle est qu'on peut paraphraser certains des arguments et des résultats du paragraphe 6, pour $\omega_0$ le foncteur identique de $\mathcal{T}$ dans $\mathcal{T}$. Dans la catégorie but le rôle essentiel est joué par les $\underline{\Hom}(X, Y)$, et le foncteur identique est à interpréter comme attachant $X$ à $X$, et comme définissant
\[
\Hom(X, Y) \to \underline{\Hom}(X, Y)
\]
pour $X, Y$ dans $\mathcal{T}$. À gauche, $\Hom(X, Y)$ est un $k$-espace vectoriel de dimension finie, interprété comme un objet de $\mathcal{T}$ (cf.~2.9 et 7.11).

\subsection{}

Soient $k$ un corps commutatif, $\mathcal{T}$ une catégorie tensorielle sur $k$, $A$ une $k$-algèbre de rang fini, $(A)_{\coh}$ la catégorie des $A$-modules à droite de type fini et $_A M$ un objet de $\mathcal{T}$ muni d'une structure de $A$-module à gauche.

Il définit un foncteur
\[
\omega_0 : (A)_{\coh} \to \mathcal{T} : E \mapsto E \otimes_A {_A M}.
\]
La catégorie $(A)_{\coh} \otimes \mathcal{T}$ est la catégorie $(\mathcal{T} \text{--} A)$ des objets de $\mathcal{T}$ munis d'une structure de $A$-module à droite (5.11). Le foncteur
\[
\omega : (\mathcal{T} \text{--} A) \to \mathcal{T} : E \mapsto E \otimes_A {_A M}
\]
est exact à droite; il correspond au foncteur
\[
(A)_{\coh} \otimes \mathcal{T} \to \mathcal{T} : (E, T) \mapsto (E \otimes_A {_A M}) \otimes T
\]
de $(A)_{\coh} \otimes \mathcal{T} \times \mathcal{T}$ dans $\mathcal{T}$.

Soit $M_A$ le dual (au sens de $\mathcal{T}$) de $_A M$. On le munit de la structure de $A$-module à droite transposée de la structure de $A$-module de $_A M$. Le morphisme d'évaluation: $M_A \otimes {_A M} \to \mathbf{1}$ se factorise par $\ev_A : M_A \otimes_A {_A M} \to \mathbf{1}$ et le morphisme $\delta : \mathbf{1} \to {_A M} \otimes M_A$ se prolonge de façon unique en un morphisme de $A$-bimodules $\delta_A : A \to {_A M} \otimes M_A = \underline{\Hom}({_A M}, {_A M})$. Soit $\omega^*$ le foncteur de $\mathcal{T}$ dans $(\mathcal{T} \text{--} A)$:
\[
\omega^* : T \mapsto T \otimes M_A.
\]
On laisse au lecteur le soin de vérifier que, comme en 4.3, $\ev_A$ et $\delta_A$ définissent des morphismes de foncteurs $\omega\omega^* \to \Id$ et $\Id \to \omega^*\omega$ faisant de $(\omega, \omega^*)$ une paire de foncteurs adjoints. Posons $L := M_A \otimes_A {_A M}$. Le foncteur $\omega\omega^*$ de $\mathcal{T}$ dans $\mathcal{T}$ est $\otimes L$, et le morphisme de foncteur $\omega\omega^* \to \omega(\omega^*\omega)\omega^* = (\omega\omega^*)(\omega\omega^*)$ est défini par un coproduit $c : L \to L \otimes L$, coassociatif et de counité $\ev_A$ : on dira que $L$ est une \emph{cogèbre} de $\mathcal{T}$.

Une \emph{représentation} de $L$ est un objet $T$ de $\mathcal{T}$ muni de $\mu : T \to T \otimes L$ rendant commutatif
\[
\begin{tikzcd}
T \arrow[r, "\mu"] \arrow[d, "\mu"'] & T \otimes L \arrow[d, "\mu \otimes \Id_L"] \\
T \otimes L \arrow[r, "\Id_T \otimes c"'] & T \otimes L \otimes L
\end{tikzcd}
\]
et tel que $T \to T \otimes L \xrightarrow{\Id \otimes \ev_A} T$ soit l'identité. La cogèbre $L$ de $\mathcal{T}$ agit fonctoriellement sur chaque $\omega(E)$.

Supposons $\omega_0$ exact et fidèle, et que $\mathcal{T}$ vérifie (2.12.1). D'après 5.14(i), le foncteur
\[
\omega_0 \otimes \Id : (A)_{\coh} \otimes \mathcal{T} \to \mathcal{T} \otimes \mathcal{T}
\]
est encore exact et fidèle. Si $\mathcal{T} \otimes \mathcal{T}$ est tensorielle, le foncteur $k$-linéaire exact à droite $T : \mathcal{T} \otimes \mathcal{T} \to \mathcal{T}$ tel que $T(X \otimes Y) = X \otimes Y$ est exact et fidèle (2.10). Le foncteur composé $\omega = T \circ (\omega_0 \otimes \Id)$ est exact et fidèle et le théorème de Barr-Beck (4.1) donne:

\begin{proposition}[8.3]
Avec les hypothèses et notations de 8.2, si $\mathcal{T}$ vérifie (2.12.1), que $\mathcal{T} \otimes \mathcal{T}$ est tensorielle et que $\omega_0 : (A)_{\coh} \to \mathcal{T}$ est exact et fidèle, le foncteur $\omega$ induit une équivalence de $(A)_{\coh} \otimes \mathcal{T}$ avec la catégorie des objets de $\mathcal{T}$ munis d'une action de $L$.
\end{proposition}

\subsection{(parallèle à 4.7)}

Soient $C$ une petite catégorie et $\omega_1, \omega_2$ deux foncteurs de $C$ dans $\mathcal{T}$. Nous avons défini en 6.12 un objet $\Lambda_k(\omega_1, \omega_2)$ de $\Ind(\mathcal{T})$. Si cela ne crée pas de confusion, nous le noterons simplement $\Lambda(\omega_1, \omega_2)$.

On peut interpréter (6.12.1) comme un morphisme
\begin{equation}\label{eq:8.4.1}
\Lambda(X) : \omega_2(X) \to \omega_1(X) \otimes \Lambda(\omega_1, \omega_2)
\end{equation}
et (6.12.1) signifie que ce morphisme est fonctoriel en $X$. La propriété universelle de $\Lambda := \Lambda(\omega_1, \omega_2)$ est que pour tout ind-objet $U$ de $\mathcal{T}$, l'application $f \mapsto f \circ \Lambda(X)$ de $\Hom(\Lambda, U)$ dans les systèmes de morphismes $u(X) : \omega_2(X) \to \omega_1(X) \otimes U$ fonctoriels en $X$, est une bijection.

Pour trois foncteurs $\omega_1, \omega_2, \omega_3$, itérant (8.4.1) on obtient
\[
\omega_3(X) \to \omega_2(X) \otimes \Lambda(\omega_2, \omega_3) \to \omega_1(X) \otimes \Lambda(\omega_1, \omega_2) \otimes \Lambda(\omega_2, \omega_3)
\]
fonctoriel en $X$, d'où
\begin{equation}\label{eq:8.4.2}
\Lambda(\omega_1, \omega_3) \to \Lambda(\omega_1, \omega_2) \otimes \Lambda(\omega_2, \omega_3).
\end{equation}
Le coproduit (8.4.2) est coassociatif en un sens évident et les morphismes d'évaluation $\omega_i(X)^\vee \otimes \omega_i(X) \to \mathbf{1}$ définissent comme en 4.7 une counité
\[
\varepsilon : \Lambda(\omega_i, \omega_i) \to \mathbf{1}.
\]
Pour $\omega$ un foncteur de $C$ dans $\mathcal{T}$, posons $\Lambda(\omega) := \Lambda(\omega, \omega)$. Le coproduit (8.4.4) de $\Lambda(\omega)$:
\begin{equation}\label{eq:8.4.3}
\Lambda(\omega) \to \Lambda(\omega) \otimes \Lambda(\omega)
\end{equation}
est coassociatif et admet une counité.

\begin{example}[8.5 (parallèle à 4.8)]
\textup{(i)} Pour $C$ réduite à un objet et à sa flèche identique, la donnée de $\omega$ est celle d'un objet $M$ de $\mathcal{T}$ et $\Lambda(\omega) \simeq M^\vee \otimes M$.

\textup{(ii)} Soient $A$ une algèbre de dimension finie sur $k$, $_A M$ un objet de $\mathcal{T}$ muni d'une structure de $A$-module à gauche, $C$ une sous-catégorie pleine de la catégorie des $A$-modules à droite de type fini, contenant $A$, et $\omega$ la restriction à $C$ du foncteur $E \mapsto E \otimes_A {_A M}$. Le morphisme $\omega(A)^\vee \otimes \omega(A) \to \Lambda(\omega)$ : $(_A M)^\vee \otimes {_A M} \to \Lambda(\omega)$ se factorise par un isomorphisme
\[
(_A M)^\vee \otimes {_A M} \xrightarrow{\sim} \Lambda(\omega).
\]
Pour $C$ réduite à $A$, c'est la définition du produit tensoriel sur $A$. Le cas général sera traité en 8.7.
\end{example}

\subsection{Fonctorialités (parallèle à 4.10)}

\textup{(i)} Soient $\eta : \mathcal{T} \to \mathcal{T}'$ un $\otimes$-foncteur exact à droite, donc exact d'après 2.10, entre catégories tensorielles. Notant encore $\eta$ son extension aux Ind-objets, on a
\begin{equation}\label{eq:8.6.1}
\eta \Lambda(\omega_1, \omega_2) \xrightarrow{\sim} \Lambda(\eta\omega_1, \eta\omega_2).
\end{equation}

\textup{(ii)} Soit $\eta : \mathcal{D} \to C$. Les morphismes $\omega_2(\eta D) \to \omega_1(\eta D) \otimes \Lambda(\omega_1, \omega_2)$ ($D$ dans $\mathcal{D}$) définissent
\begin{equation}\label{eq:8.6.2}
\Lambda(\omega_1 \circ \eta, \omega_2 \circ \eta) \to \Lambda(\omega_1, \omega_2).
\end{equation}

\textup{(iii)} Si $C$ est limite inductive filtrante de catégories $C_i$ et que $\mathcal{T}_i$ est le foncteur naturel de $C_i$ dans $C$, les morphismes (8.6.2) induisent un isomorphisme
\begin{equation}\label{eq:8.6.3}
\varinjlim \Lambda(\omega_1 \circ \mathcal{T}_i, \omega_2 \circ \mathcal{T}_i) \xrightarrow{\sim} \Lambda(\omega_1, \omega_2).
\end{equation}

\textup{(iv)} Soient $(C_i)_{i \in I}$ une famille finie de catégories, $\omega_i^j$ de source $C_i$ ($i \in I$, $j = 1, 2$), $C$ le produit des $C_i$ et $\otimes \omega_i^j$ le foncteur de $C$ dans $\mathcal{T}$ produit tensoriel des $\omega_i^j$ : $(\otimes \omega_i^j)((C_i)) = \otimes \omega_i^j(C_i)$. On a
\begin{equation}\label{eq:8.6.4}
\otimes_i \Lambda(\omega_i^1, \omega_i^2) \xrightarrow{\sim} \Lambda(\otimes \omega_i^1, \otimes \omega_i^2).
\end{equation}

\subsection{Vérification de 8.5\textup{(ii)} (parallèle à 4.11)}

Soit $L := M_A \otimes_A {_A M}$. Par 8.2, $L$ coagit sur les $\omega(E)$. La propriété universelle de $\Lambda(\omega)$ fournit $\Lambda(\omega) \to L$. Appliquant 8.6(ii) à la sous-catégorie pleine $\mathcal{D}$ de $C$ réduite à $A$, on obtient $L \to \Lambda(\omega)$ (8.5(ii) pour $\mathcal{D}$) et on vérifie que le composé $L \to \Lambda(\omega) \to L$ est l'identité. Prouvons que $L \to \Lambda(\omega)$ est un épimorphisme, i.e.~que pour tout $E$ dans $C$, l'image de $\omega(E)^\vee \otimes \omega(E)$ dans $\Lambda(\omega)$ est contenue dans celle de $\omega(A)^\vee \otimes \omega(A)$. Soit $u$ un épimorphisme $A^n \to E$, de composantes $u_i$. Les diagrammes
\[
\begin{tikzcd}
\omega(E)^\vee \otimes \omega(A) \arrow[r] \arrow[d] & \omega(A)^\vee \otimes \omega(A) \arrow[d] \\
\omega(E)^\vee \otimes \omega(E) \arrow[r] & \Lambda(\omega)
\end{tikzcd}
\]
sont commutatifs et on conclut en observant que la somme
\[
\omega(A)^{\vee n} \otimes \omega(A) \to \omega(E)^\vee \otimes \omega(E)
\]
est un épimorphisme.

\subsection{}

Soient $\mathcal{T}$ une catégorie tensorielle sur $k$, $\mathcal{A}_1$ et $\mathcal{A}_2$ deux catégories abéliennes $k$-linéaires vérifiant (2.12.1), $\omega_i : \mathcal{A}_i \to \mathcal{T}$ $k$-linéaire et exact à droite, et $\omega : \mathcal{A}_1 \otimes \mathcal{A}_2 \to \mathcal{T}$ $k$-linéaire et exact à droite défini par
\[
\omega(X_1 \otimes X_2) = \omega_1(X_1) \otimes \omega_2(X_2).
\]
Appliquant 8.6(ii) à $\otimes : \mathcal{A}_1 \times \mathcal{A}_2 \to \mathcal{A}_1 \otimes \mathcal{A}_2$, on obtient
\begin{equation}\label{eq:8.8.1}
\Lambda(\omega_1) \otimes \Lambda(\omega_2) \xrightarrow{\sim} \Lambda(\omega_1(X_1) \otimes \omega_2(X_2)) \to \Lambda(\omega).
\end{equation}
Ce morphisme est un isomorphisme. Pour le vérifier, on écrit $\mathcal{A}_i = \varinjlim \langle X_i^\alpha \rangle$ ($X_i$ dans $\mathcal{A}_i$) et on se ramène par 5.1 et 8.6(iii) à supposer $\mathcal{A}_i = \langle X_i \rangle$, donc que $\mathcal{A}_i$ est de la forme $(A_i)_{\coh}$ avec $A_i$ de dimension finie sur $k$. Le foncteur $\omega_i$ s'écrit $E \mapsto E \otimes M_i$, avec $M_i$ un objet de $\mathcal{T}$ muni d'une structure de $A_i$-module. On a $\mathcal{A}_1 \otimes \mathcal{A}_2 = (A_1 \otimes A_2)_{\coh}$, $\omega$ est $\otimes_{A_1 \otimes A_2}(M_1 \otimes M_2)$ et le morphisme s'identifie à l'isomorphisme évident
\[
(M_1^\vee \otimes_{A_1} M_1) \otimes (M_2^\vee \otimes_{A_2} M_2) \xrightarrow{\sim} (M_1 \otimes M_2)^\vee \otimes_{A_1 \otimes A_2} (M_1 \otimes M_2).
\]

\subsection{(parallèle à 6.3)}

Soient $\mathcal{A}_i$ ($i = 1, 2, 3$) trois catégories et $\omega_j^i$ (resp.~$\omega_k^i$) un foncteur de $\mathcal{A}_i$ dans $\mathcal{T}$. Supposons donné un foncteur $\otimes : \mathcal{A}_1 \times \mathcal{A}_2 \to \mathcal{A}_3$ et, pour $j = 1, 2$ un isomorphisme de foncteurs
\[
\omega_j^1(X_1) \otimes \omega_j^2(X_2) \xrightarrow{\sim} \omega_j^3(X_1 \otimes X_2).
\]
Appliquant 8.6(ii)(iv), on obtient un produit
\begin{equation}\label{eq:8.9.1}
\Lambda(\omega_1^1, \omega_2^1) \otimes \Lambda(\omega_1^2, \omega_2^2) \to \Lambda(\omega_1^3, \omega_2^3).
\end{equation}
caractérisé par l'exigence que, quels que soient $X_1$ et $X_2$ dans $\mathcal{A}_1$ et $\mathcal{A}_2$, le morphisme (8.4.1)
\[
\omega_2^3(X_1 \otimes X_2) \to \omega_1^3(X_1 \otimes X_2) \otimes \Lambda(\omega_1^3, \omega_2^3)
\]
se déduit des morphismes analogues pour les $\omega_2^i(X_i)$ et $\omega_1^i(X_i)$ ($i = 1, 2$) par (8.9.1).

Soit $\mathcal{A}$ une catégorie à produit tensoriel vérifiant (1.1.1) et $\omega_1, \omega_2$ deux $\otimes$-foncteurs de $\mathcal{A}$ dans $\mathcal{T}$. Faisons $\mathcal{A}_i = \mathcal{A}$ et $\omega_j^i = \omega_j$ ($i = 1, 2, 3$). Le produit (8.9.1) devient un produit
\begin{equation}\label{eq:8.9.2}
\Lambda(\omega_1, \omega_2) \otimes \Lambda(\omega_1, \omega_2) \to \Lambda(\omega_1, \omega_2).
\end{equation}

La même preuve qu'en 6.4 donne
\begin{proposition}[8.10]
Le produit (8.9.2) est associatif, commutatif et à unité.
\end{proposition}

Pour $T = \Sp(B)$ un schéma affine en $\mathcal{T}$ et $\omega$ un $\otimes$-foncteur de $\mathcal{A}$ dans les modules sur $T$, notons $\omega_T$ le $\otimes$-foncteur $A \mapsto \omega(A) \otimes B$ de $\mathcal{A}$ dans les modules sur $T$.
La preuve de 6.6 donne
\begin{proposition}[8.11]
Sous les hypothèses de 8.9.2,

\textup{(i)} Le $T$-schéma $\Sp \Lambda(\omega_1, \omega_2)$ représente le foncteur $\underline{\Hom}^\otimes(\omega_1, \omega_2)$ qui au $T$-schéma affine $T$ attache l'ensemble des morphismes de $\otimes$-foncteurs de $\omega_{1T}$ dans $\omega_{2T}$.

\textup{(ii)} Si $\mathcal{A}$ vérifie (1.1.1) et (1.1.2), tout morphisme $\omega_{2T} \to \omega_{1T}$ comme en (i) est un isomorphisme: $\underline{\Hom}^\otimes(\omega_1, \omega_2)$ coïncide avec $\underline{\Isom}^\otimes(\omega_1, \omega_2)$.

\textup{(iii)} Étant donnés trois $\otimes$-foncteurs $\omega_1, \omega_2, \omega_3$ la composition
\[
\underline{\Hom}^\otimes(\omega_1, \omega_2) \times \underline{\Hom}^\otimes(\omega_2, \omega_3) \to \underline{\Hom}^\otimes(\omega_1, \omega_3)
\]
est définie par (8.4.2).
\end{proposition}

\subsection{}

Prenons pour $\mathcal{A}$ la catégorie $\mathcal{T}$ et pour foncteurs $\omega_1$ et $\omega_2$ le foncteur identique. Posons $\Lambda = \Lambda(\Id_\mathcal{T})$. D'après 8.11, $\Sp(\Lambda)$ est le $\mathcal{T}$-groupe qui représente le foncteur qui à $T = \Sp(B)$ attache le groupe des automorphismes de $\otimes$-foncteur du foncteur $X \mapsto X_T = X \otimes B$ de $\mathcal{T}$ dans les modules sur $T$.

\begin{definition}[8.13]
Le $\mathcal{T}$-groupes $\Sp(\Lambda)$ construit en 8.12 est le \emph{groupe fondamental} de $\mathcal{T}$. On le note $\pi(\mathcal{T})$.
\end{definition}

Si $\omega$ est un foncteur fibre de $\mathcal{T}$ sur un schéma $S$ sur $k$, $\omega(\pi(\mathcal{T}))$ est un schéma en groupes affines sur $S$. L'action donnée par (8.4.1) de $\pi(\mathcal{T})$ sur chaque $X$ de $\mathcal{T}$ définit une action de $\omega(\pi(\mathcal{T}))$ sur les $\omega(X)$. Cette action est fonctorielle et compatible au produit tensoriel. Comparant 4.7 et 8.2--8.4, on trouve que
\begin{equation}\label{eq:8.13.1}
\omega(\pi(\mathcal{T})) \xrightarrow{\sim} \underline{\Aut}_k(\omega).
\end{equation}

A.~Grothendieck compare cette situation à la suivante: pour $X$ un espace topologique localement connexe et localement simplement connexe, les groupes fondamentaux $\pi_1(X, x)$ forment un système local sur $X$, i.e.~un objet en groupes de la catégorie des faisceaux localement constants sur $X$. Ce faisceau de groupes agit sur tout faisceau localement constant. Sa fibre $\pi_1(X, x)$ en $x$ est le groupe des automorphismes du foncteur ``fibre en $x$''
\[
\text{(faisceaux localement constants)} \to \text{(Ens)}.
\]

\begin{example}[8.14]
\textup{(i)} Pour $\mathcal{T}$ une catégorie tannakienne, la donnée d'un $\mathcal{T}$-groupe $G$ équivaut par $G \mapsto (\omega(G))$ à la donnée, pour tout foncteur fibre $\omega$ sur tout schéma $S$ sur $k$, d'un schéma en groupes affines $G_\omega$ sur $S$, de formation compatible à tout changement de base $S'/S$. Le groupe fondamental $\pi(\mathcal{T})$ est donné par
\[
\omega \mapsto \underline{\Aut}_k(\omega).
\]

\textup{(ii)} Cas particulier: si $\mathcal{T} = \Rep(G')$, un $\mathcal{T}$-groupe s'identifie à un schéma en groupes affines sur $k$, muni d'une action de $G'$. Le schéma en groupe $G'$, muni de son action sur lui-même par automorphismes intérieurs, correspond à $\pi(\Rep(G'))$.

\textup{(iii)} Pour $G'$ le groupe trivial, on trouve que $\pi(\Vect k)$ est le groupe trivial.

\textup{(iv)} Prenons pour $\mathcal{T}$ la catégorie des super-espaces vectoriels de dimension finie sur $k$ (1.4). Regardons le $k$-schéma en groupes $\mu_2$ des racines carrées de l'unité comme un $\mathcal{T}$-groupe (cf.~7.11). Il agit sur chaque super-espace vectoriel par l'action ``parité'': $(\varepsilon, a) \mapsto \varepsilon^{\deg a} a$ pour $a$ homogène, et on vérifie que
\[
\pi(\mathcal{T}) = \mu_2.
\]
\end{example}

\subsection{Fonctorialité}

Soit $\eta : \mathcal{T}_1 \to \mathcal{T}_2$ un $\otimes$-foncteur exact $k$-linéaire entre catégories tensorielles sur $k$. Faisant $\mathcal{T} = \eta$ dans 8.6(i), on obtient
\begin{equation}\label{eq:8.15.1}
\Sp(\Lambda(\eta)) = \eta(\pi(\mathcal{T}_1))
\end{equation}
et le $\mathcal{T}_2$-schéma en groupes affines $\eta(\pi(\mathcal{T}_1))$ représente le foncteur $\Aut^\otimes(\eta) : T = \Sp(B) \mapsto$ automorphismes du $\otimes$-foncteur $X \mapsto \eta_{\mathcal{T}_2}(X) = \eta(X) \otimes B$ de $\mathcal{T}_1$ dans les modules sur $T$. Faisant $\mathcal{T} = \mathcal{T}_2$ dans 8.6(ii), on obtient un morphisme
\begin{equation}\label{eq:8.15.2}
\pi(\mathcal{T}_2) \to \eta(\pi(\mathcal{T}_1)).
\end{equation}
Le morphisme de foncteurs correspondant attache à un automorphisme du foncteur $X \mapsto X_T = X \otimes B$ l'automorphisme qu'il induit de son composé avec $\eta$.

\subsection{Produit}

Soient $\mathcal{T}_1$ et $\mathcal{T}_2$ deux catégories tensorielles sur $k$. On suppose que $\mathcal{T}_1$ et $\mathcal{T}_2$ vérifient (2.12.1) et que $\mathcal{T} := \mathcal{T}_1 \otimes \mathcal{T}_2$ est encore tensorielle sur $k$. On identifie $\mathcal{T}_1$ à une sous-catégorie pleine de $\mathcal{T}$ stable par sous-quotients par
\[
\inj_1 : \mathcal{T}_1 = \mathcal{T}_1 \otimes \Vect(k) \hookrightarrow \mathcal{T}_1 \otimes \mathcal{T}_2
\]
\[
X \mapsto X \otimes \mathbf{1}
\]
(appliquer 2.9 et 5.14(i)). De même pour $\mathcal{T}_2$. Le produit tensoriel
\[
\otimes_k : \mathcal{T} \times \mathcal{T} \to \mathcal{T} \otimes \mathcal{T}
\]
est le composé du produit tensoriel dans $\mathcal{T} \otimes \mathcal{T}$ et de $(\inj_1, \inj_2)$:
\[
X_1 \otimes X_2 = \inj_1(X_1) \otimes \inj_2(X_2).
\]
Par 8.8 appliqué au foncteur identique de $\mathcal{T} \otimes \mathcal{T}$ dans $\mathcal{T} \otimes \mathcal{T}$, $\pi(\mathcal{T} \otimes \mathcal{T})$ est le $\Sp$ de $\Lambda(\otimes_k : \mathcal{T} \times \mathcal{T} \to \mathcal{T} \otimes \mathcal{T})$. Par 8.6(iv), ce $\Lambda$ est le produit tensoriel des $\Lambda(\inj_i : \mathcal{T} \to \mathcal{T} \times \mathcal{T})$ qui, par 8.6(i) pour $\eta = \inj_i$ et $\omega = \Id_\mathcal{T}$ coïncident avec $\inj_i \Lambda(\Id_\mathcal{T})$. On conclut que les morphismes (8.15.2) (pour $\eta = \inj_i$)
\[
\pi(\mathcal{T}_1 \otimes \mathcal{T}_2) \to \inj_i \pi(\mathcal{T}_i)
\]
définissent un isomorphisme
\begin{equation}\label{eq:8.16.1}
\pi(\mathcal{T}_1 \otimes \mathcal{T}_2) \xrightarrow{\sim} \inj_1(\pi(\mathcal{T}_1)) \times \inj_2(\pi(\mathcal{T}_2)).
\end{equation}

\begin{theorem}[8.17]
Soient $\mathcal{T}_1$ et $\mathcal{T}_2$ deux catégories tensorielles sur $k$ vérifiant (2.12.1) et $\eta$ un $\otimes$-foncteur exact de $\mathcal{T}_1$ dans $\mathcal{T}_2$. On suppose que $\mathcal{T}_2 \otimes \mathcal{T}_2$ est encore tensorielle sur $k$ (cf.~8.1). Le $\otimes$-foncteur induit par $\eta$, de $\mathcal{T}_1$ dans la catégorie des objets de $\mathcal{T}_2$ munis d'une action de $\eta(\pi(\mathcal{T}_1))$ tels que l'action naturelle de $\pi(\mathcal{T}_2)$ s'en déduise par (8.15.2), est une équivalence de catégories.
\end{theorem}

La preuve sera donnée en 8.26.

Prenant pour $\mathcal{T}_1$ la catégorie $\Vect(k)$ des multiples de $\mathbf{1}$, on obtient par 8.14(iii):

\begin{corollary}[8.18]
Soit $\mathcal{T}$ une catégorie tensorielle sur $k$. On suppose que $\mathcal{T}$ vérifie (2.12.1) et que $\mathcal{T} \otimes \mathcal{T}$ est encore tensorielle sur $k$. Alors, pour que $X$ dans $\mathcal{T}$ soit une somme de copies de $\mathbf{1}$, il faut et il suffit que $\pi(\mathcal{T})$ agisse trivialement sur $X$.
\end{corollary}

Si on prend pour $\mathcal{T}_2$ la catégorie $\Vect(k)$, on retrouve le fait qu'un foncteur fibre $\omega$ de $\mathcal{T}$ sur $k$ induit une équivalence $\mathcal{T} \to \Rep_k(\underline{\Aut}_k(\omega))$, puisque $\omega(\pi(\mathcal{T})) = \underline{\Aut}_k(\omega)$.

\subsection{}

Un autre cas intéressant est celui où $\mathcal{T}_2$ est la catégorie 1.4 des super-espaces vectoriels de dimension finie. On a $\pi(\mathcal{T}_2) = \mu_2$ (8.14(iv)). Soit $G$ est un super groupe affine, i.e.~un $\mathcal{T}_2$-groupes, muni de $p : \mu_2 \to G$ tel que l'action par automorphismes intérieurs de $\mu_2$ sur l'algèbre affine de $G$ soit l'action ``parité'' (8.14(iv)). La catégorie des super-espaces vectoriels $V$ de dimension finie, munis d'une action $\rho$ de $G$ telle que $\rho \circ p$ soit l'action ``parité'' est une catégorie tensorielle $\Rep(G, p)$ sur $k$. D'après 8.17, toute catégorie tensorielle $\mathcal{T}$ sur $k$, muni d'un $\otimes$-foncteur exact
\[
\omega : \mathcal{T} \to \text{(super-espaces vectoriels)}
\]
est ainsi obtenue: on a une équivalence
\[
\mathcal{T} \xrightarrow{\sim} \Rep(\omega(\pi(\mathcal{T})), p \text{ donné par (8.15.2)}).
\]

\subsection{}

Soient $k$ un corps, $\mathcal{T}$ une catégorie tensorielle sur $k$, $\mathcal{A}$ une catégorie abélienne vérifiant (2.12.1) et $\omega_0$ un foncteur $k$-linéaire exact à droite de $\mathcal{A}$ dans $\mathcal{T}$. La catégorie $\mathcal{A}$ est limite inductive filtrante de ses sous-catégories $\langle X \rangle$, et chacune est équivalente à une catégorie $(A)_{\coh}$ pour $A$ une $k$-algèbre de rang fini. Par (5.1) et 5.11, $\mathcal{A} \otimes \mathcal{T}$ existe:
\[
\mathcal{A} \otimes \mathcal{T} = \varinjlim \langle X \rangle \otimes \mathcal{T}.
\]
Le foncteur $\mathcal{A} \times \mathcal{T} \to \mathcal{T} : (\omega_0(A), T) \mapsto \omega_0(A) \otimes T$ est exact à droite et se factorise donc uniquement par
\[
\omega : \mathcal{A} \otimes \mathcal{T} \to \mathcal{T}.
\]
On a
\[
\Lambda(\omega_0) = \varinjlim \Lambda(\omega_0 | \langle X \rangle).
\]
Par définition des Ind-objets, on a pour $T$ dans $\mathcal{T}$
\[
\Hom(T, T \otimes \Lambda(\omega_0)) = \varinjlim \Hom(T, T \otimes \Lambda(\omega_0 | \langle X \rangle))
\]
et la catégorie des objets de $\mathcal{T}$ munis d'une coaction de $\Lambda(\omega_0)$ est la limite inductive des catégories analogues pour les $\Lambda(\omega_0 | \langle X \rangle)$. Appliquant 8.3 et 8.5(ii), on obtient:

\begin{proposition}[8.21]
Avec les hypothèses et notations de 8.19, si $\mathcal{T} \otimes \mathcal{T}$ est tensorielle et que $\omega_0 : \mathcal{A} \to \mathcal{T}$ est exact et fidèle, le foncteur $\omega$ induit une équivalence de $\mathcal{A} \otimes \mathcal{T}$ avec la catégorie des objets de $\mathcal{T}$ munis d'une coaction de $\Lambda(\omega_0)$.
\end{proposition}

En particulier, par (8.15.1)
\begin{proposition}[8.22]
Soit $\eta : \mathcal{T}_1 \to \mathcal{T}_2$ un $\otimes$-foncteur exact $k$-linéaire entre catégories tensorielles sur $k$. Si $\mathcal{T}_1$ vérifie (2.12.1) et que $\mathcal{T} \otimes \mathcal{T}$ est tensorielle, le foncteur $\omega : \mathcal{T}_1 \otimes \mathcal{T} \to \mathcal{T}$ tel que $\omega(T_1 \otimes T) = \eta(T_1) \otimes T$ induit une équivalence de $\mathcal{T}_1 \otimes \mathcal{T}$ avec la catégorie des objets de $\mathcal{T}$ munis d'une action de $\eta(\pi(\mathcal{T}_1))$.
\end{proposition}

Supposons que $\mathcal{T}$ vérifie (2.12.1) et faisons $\mathcal{T}_1 = \mathcal{T}$. Prenons pour foncteur $\omega_0$ le foncteur identique de $\mathcal{T}$ dans $\mathcal{T}$. On obtient
\begin{corollary}[8.23]
Supposons que $\mathcal{T}$ vérifie (2.12.1) et que $\mathcal{T} \otimes \mathcal{T}$ est tensorielle. Le foncteur $\omega : \mathcal{T} \otimes \mathcal{T} \to \mathcal{T}$ tel que $\omega(X \otimes Y) = X \otimes Y$ induit une équivalence de $\mathcal{T} \otimes \mathcal{T}$ avec la catégorie des objets de $\mathcal{T}$ munis d'une action de $\pi(\mathcal{T})$; l'action de $\pi(\mathcal{T})$ sur $\omega(X \otimes Y) = X \otimes Y$ est celle déduite de son action sur $X$.
\end{corollary}

\subsection{Preuve de 8.18}

La sous-catégorie $\Vect(k) \subset \mathcal{T}$ (2.9) est stable par sous-quotient et
\[
\Vect(k) \otimes \mathcal{T} \hookrightarrow \mathcal{T} \otimes \mathcal{T}
\]
est donc pleinement fidèle d'image essentielle stable par sous-quotients (5.14(ii)). Le groupe $\pi(\Vect(k))$ est trivial et comparant 8.22 pour $\mathcal{T}_1 = \Vect(k)$ et $\mathcal{T}$, on trouve que l'équivalence 8.23 de $\mathcal{T} \otimes \mathcal{T}$ avec les objets de $\mathcal{T}$ à action de $\pi(\mathcal{T})$ identifie $\Vect(k) \otimes \mathcal{T}$ à la sous-catégorie des objets de $\mathcal{T}$ à action triviale. En particulier, pour $X$ dans $\mathcal{T}$, $X \otimes \mathbf{1}$ est dans $\Vect(k) \otimes \mathcal{T}$ si et seulement si l'action naturelle de $\pi(\mathcal{T})$ sur $X$ est triviale. D'après 5.14(ii), pour que $X \otimes \mathbf{1}$ soit dans $\Vect(k) \otimes \mathcal{T}$, il faut et il suffit que $X$ soit dans $\Vect(k)$ et 8.18 en résulte.

\begin{corollary}[8.25]
Soient $\mathcal{T}_1, \mathcal{T}_2$ deux catégories tensorielles sur $k$ vérifiant (2.12.1) et $\mathcal{T} = \mathcal{T}_1 \otimes \mathcal{T}_2$. On suppose $\mathcal{T}_2 \otimes \mathcal{T}_2$ et $\mathcal{T}$ tensorielles. Le foncteur $\inj_1 : \mathcal{T}_1 \to \mathcal{T}_1 \otimes \mathcal{T}_2 : X \mapsto X \otimes \mathbf{1}$ est une équivalence de $\mathcal{T}_1$ avec la catégorie des objets de $\mathcal{T}_1 \otimes \mathcal{T}_2$ sur lesquels $\inj_2(\pi(\mathcal{T}_2))$ agit trivialement.
\end{corollary}

Puisque $\pi(\mathcal{T}) \xrightarrow{\sim} \inj_1 \pi(\mathcal{T}_1) \times \inj_2 \pi(\mathcal{T}_2)$, on a une inclusion $\inj_1 \pi(\mathcal{T}_1) \hookrightarrow \pi(\mathcal{T})$; via cette inclusion, $\inj_1 \pi(\mathcal{T}_1)$ agit sur les objets de $\mathcal{T}_1 \otimes \mathcal{T}_2$. C'est de cette action qu'il est question dans le corollaire.

\begin{proof}
Si un $\mathcal{T}$-groupe $G$ agit sur un objet $T$ de $\mathcal{T}$, $G = \Sp(B)$, l'action de $G$ est un morphisme $\rho : T \to T \otimes B$. L'unité de $B$ : $\mathbf{1} \to B$ définit une autre flèche $\Id \otimes 1 : T \to T \otimes B$. Soit $T^G$ le noyau de la double flèche $(\rho, \Id \otimes 1)$. C'est un Ind-objet de $\mathcal{T}$, et un sous-objet de $T$. Parce que $\mathcal{T}$ est noethérienne (5.13.2), tout Ind-objet de $\mathcal{T}$ est réunion croissante de ses sous-objets dans $\mathcal{T}$. Un Ind-objet, qui est sous-objet d'un objet de $\mathcal{T}$, est donc dans $\mathcal{T}$ et $T^G$ est dans $\mathcal{T}$.

Appliquons 8.18. Le foncteur $\mathcal{T}_2 \to \mathcal{T}_2 : Y \mapsto Y^{\pi(\mathcal{T}_2)}$ se factorise par $\Vect(k)$ et est exact à gauche. Il définit un foncteur exact à gauche $\mathcal{T}_1 \otimes \mathcal{T}_2 \to \mathcal{T}_1 \otimes \mathcal{T}_2$ se factorisant par $\mathcal{T}_1 \otimes \Vect(k) \simeq \mathcal{T}_1$. Ce foncteur coïncide avec $T \mapsto T^{\inj_2(\pi(\mathcal{T}_2))}$: il suffit de le vérifier après composition après $\mathcal{T}_1 \times \mathcal{T}_2 \to \mathcal{T}_1 \otimes \mathcal{T}_2$, et que
\[
(X_1 \otimes X_2)^{\inj_2(\pi(\mathcal{T}_2))} = X_1 \otimes X_2^{\pi(\mathcal{T}_2)}
\]
résulte de l'exactitude de $\otimes_k$. En particulier, si $T$ est fixe par $\inj_2(\pi(\mathcal{T}_2))$, $T$ est dans l'image de $\mathcal{T}_1$. La réciproque est claire.
\end{proof}

\subsection{Preuve de 8.17}

On sait déjà (8.22) que
\[
\eta \otimes \Id : \mathcal{T}_1 \otimes \mathcal{T}_2 \to \text{($\eta(\pi(\mathcal{T}_1))$ -- objets de $\mathcal{T}_2$)}
\]
est une équivalence. En particulier, $\mathcal{T}_1 \otimes \mathcal{T}_2$ est tensorielle. Soient $X$ dans $\mathcal{T}_1 \otimes \mathcal{T}_2$. Il est muni d'une action $(\rho_1, \rho_2)$ de $\inj_1 \pi(\mathcal{T}_1) \times \inj_2 \pi(\mathcal{T}_2)$. Son image $(\eta \otimes \Id)(X)$ dans $\mathcal{T}_2$ est munie d'une action $\rho$ de $\pi(\mathcal{T}_2)$. Cette action est le produit des actions --- qui commutent --- suivantes: $\rho'_1$, déduit par (8.15.2) $\pi(\mathcal{T}_2) \to \eta \pi(\mathcal{T}_1)$ de l'action de $\inj_1 \pi(\mathcal{T}_1)$, et $\rho'_2$, déduit de l'action $\rho_2$. Tout objet de $\mathcal{T}_1 \otimes \mathcal{T}$ étant quotient d'un $\mathcal{T}_1 \otimes \mathcal{T}_2$, il suffit de le vérifier dans le cas --- laissé au lecteur --- où $X = T_1 \otimes T_2$.

Le foncteur $\eta \otimes \Id$ transforme donc l'action de $\inj_2 \pi(\mathcal{T}_2)$ sur $X$ en la différence entre deux actions de $(\eta \otimes \Id) \circ \inj_2 \pi(\mathcal{T}_2) = \pi(\mathcal{T}_2)$ sur $(\eta \otimes \Id)(X)$. L'action naturelle, et l'action via (8.15.2): $\pi(\mathcal{T}_2) \to \eta \pi(\mathcal{T}_1)$. Pour que l'action de $\inj_2 \pi(\mathcal{T}_2)$ sur $X$ soit triviale (i.e.~que $X$ soit dans $\mathcal{T}_1$) il faut et suffit que son image par $\eta \otimes \Id$ le soit, et 8.17 résulte de 8.25.

\section{Corps différentiels}

\subsection{}

Soit $(K, \partial)$ un corps différentiel: un corps commutatif de caractéristique 0, muni d'une dérivation $\partial : K \to K$. Le corps des constantes $k$ est $k := \ker(\partial)$. Une connexion $\nabla$ sur un espace vectoriel $V$ sur $K$, supposé de dimension finie, est $\nabla : V \to V$, additif et vérifiant l'identité de Leibniz
\[
\nabla(\lambda v) = \partial\lambda \cdot v + \lambda \nabla v.
\]
Le produit tensoriel d'espaces vectoriels à connexion est défini par
\[
\nabla(v_1 \otimes v_2) = (\nabla v_1) \otimes v_2 + v_1 \otimes (\nabla v_2).
\]
Pour ce produit tensoriel, et ses données évidentes d'associativité et de commutativité, la catégorie $\mathcal{T}$ des espaces vectoriels de dimension finie sur $K$ munis d'une connexion est une catégorie tensorielle sur $k$. L'objet unité est $(K, \partial)$ et $k = \End(K, \partial)$. Le foncteur $\omega$ de $\mathcal{T}$ dans les espaces vectoriels sur $K$ : ``oubli de $\nabla$'' est un foncteur fibre. La catégorie $\mathcal{T}$ est donc tannakienne.

Soit $(X, \nabla)$ dans $\mathcal{T}$. Comme en 6.16, on notera $\langle X \rangle^\otimes$ la sous-catégorie pleine de $\mathcal{T}$ d'objets les sous-quotients de sommes de $X^{\otimes n} \otimes X^{\vee \otimes m}$. C'est encore une catégorie tannakienne.

\subsection{}

Supposons que $\omega_0 : \langle X \rangle^\otimes \to \Vect(k)$ soit un foncteur fibre. Si $k$ est algébriquement clos, un tel foncteur fibre existe (6.20) et est unique à isomorphisme (non unique) près car tout $\underline{\Aut}_k(\omega)$-torseur sur $k$ est trivial.

Soient $G := \underline{\Aut}_k(\omega_0) \subset \GL(\omega_0(X))$, $G_K$ déduit de $G$ par extension des scalaires à $K$ et $P$ le $G_K$-torseur $\underline{\Isom}_k(\omega_0 \otimes k, \omega | \langle X \rangle^\otimes)$. C'est un sous-schéma du $\GL(\omega_0(X))$-torseur des isomorphismes $K$-linéaires $\omega_0(X) \otimes K \xrightarrow{\sim} X$.

Ce torseur est muni d'une connexion, redonnant la connexion de $(\omega_0 \otimes_K)^P \in \Ob \langle X \rangle^\otimes$ pour $W$ une représentation de $\Aut(\omega_0)$. Pour le voir, le plus simple est d'interpréter la dérivation $\partial$ comme un automorphisme infinitésimal: pour $K[\varepsilon]$ les nombres duaux ($\varepsilon^2 = 0$), $\partial$ définit l'automorphisme $k[\varepsilon]$-linéaire de $K[\varepsilon]$ : $\exp(\varepsilon\partial) : \lambda + \mu \partial \mapsto \lambda + (\partial\lambda + \mu)\varepsilon$. Une connexion sur un objet sur $K$ (espace vectoriel, torseur \ldots) s'interprète comme un automorphisme $\exp(\varepsilon\partial)$-linéaire de l'objet sur $K[\varepsilon]$ s'en déduisant par extension des scalaires. La connexion de chaque $\omega(W)$ ($W$ dans $\langle X \rangle^\otimes$) définit une connexion sur $\underline{\Isom}_k | \langle X \rangle^\otimes$ et, par transport de structure, sur $P$.

La connexion de $P$ peut aussi être vue comme une dérivation de $\Gamma(P, \mathcal{O})$. Le schéma $P$ est un sous-schéma du schéma vectoriel défini par le $K$-espace vectoriel $\Hom_K(\omega_0(X) \otimes K, X)$. La connexion de $X$ (et la connexion de $\omega_0(X) \otimes K$ avec les $a \in \omega_0(X)$ horizontaux) définissent sur cet espace vectoriel une connexion qui induit celle de $P$. Sur l'espace $\Gamma(P, \mathcal{O}) \otimes X$ des morphismes de $P$ dans $X$, on définit $\nabla$ par $\nabla(f \otimes x) = \nabla f \otimes x + f \otimes \nabla x$. Pour $a \in \omega_0(X)$, le morphisme $p \mapsto p(a)$ est horizontal.

\begin{proposition}[9.3]
Avec les notations précédentes,

\textup{(i)} Le schéma $P$ sur $K$ est lisse et connexe.

\textup{(ii)} Soit $k(P)$ le corps des fonctions rationnelles sur $P$. Le corps des constantes de ce corps différentiel est $k$.
\end{proposition}

\begin{proof}
Le schéma $P$ est lisse car c'est un torseur sous le groupe algébrique $G_K$ sur $K$.

Le foncteur $W \mapsto (W \otimes K)^P$ est une équivalence de catégories $\Rep(G) \xrightarrow{\sim} \langle X \rangle^\otimes$. En particulier, si $(V, \nabla) = (W \otimes K)^P$, les sous-espaces vectoriels de $V$ stable sous $\nabla$ sont les $(W' \otimes K)^P$ pour $W' \subset W$ $G$-stable. Le $k$-vectoriel $V^\nabla$ des vecteurs horizontaux: $\nabla v = 0$ s'identifie à $\Hom_\mathcal{T}(\mathbf{1}, V)$ (par $f \mapsto f(1)$) et donc à $W^G$. Par passage à la limite inductive, les même résultats valent pour une représentation de $G$ non nécessairement de dimension finie. À une telle représentation correspond un espace vectoriel à connexion, réunion filtrante de sous-espaces de dimension finie stables par $\nabla$ et dans $\langle X \rangle^\otimes$.

Prenons pour $W$ la représentation régulière de $G$ : $\Gamma(G, \mathcal{O})$ sur lequel $G$ agit par translations à gauche: $f(x) \mapsto f(g^{-1}x)$. Le $\nabla$ obtenu est $\Gamma(P, \mathcal{O})$, muni de $\nabla$. Sur $G$, les seules fonctions invariantes par translation sont les constantes, et les seuls sous-schémas fermés invariants par translations sont $\emptyset$ et $G$. On obtient
\begin{equation}\label{eq:9.3.1}
k \xrightarrow{\sim} \Gamma(P, \mathcal{O})^\nabla;
\end{equation}
\begin{equation}\label{eq:9.3.2}
\text{un idéal } I \subset \Gamma(P, \mathcal{O}) \text{, stable par } \nabla \text{, est trivial: } I = \Gamma(P, \mathcal{O}) \text{ ou } 0.
\end{equation}

Si $P$ n'était pas connexe, il y aurait dans $\Gamma(P, \mathcal{O})$ des idempotents non triviaux --- nécessairement horizontaux --- et ceci contredirait (9.3.1). Ceci prouve (i).

Prouvons (ii). On ne peut pas directement appliquer le même argument. (9.3.1) $\not\simeq k(P)^\nabla$, car $k(P)$ n'est pas une représentation du groupe algébrique $G$, seulement de $G(k)$. Soit $f \in k(P)$ avec $\nabla f = 0$. Soit $I \subset \Gamma(P, \mathcal{O})$ l'idéal des $g$ tels que $gf \in \Gamma(P, \mathcal{O})$. Si $gf \in \Gamma(P, \mathcal{O})$, on a $\nabla(gf) = \nabla(g) \cdot f \in \Gamma(P, \mathcal{O})$ et l'idéal $I$ est donc stable par $\nabla$. Par (9.3.2), on a $I = \Gamma(P, \mathcal{O})$ et donc $f \in \Gamma(P, \mathcal{O})$. On conclut par (9.3.1).
\end{proof}

\subsection{}

Une \emph{extension de Picard-Vessiot} de $K$, pour $(X, \nabla)$, est un corps différentiel $(L, \partial)$ extension de $(K, \partial)$ tel que
\begin{enumerate}
\item[(a)] $X \otimes_K L$ est engendré sur $L$ par ses vecteurs horizontaux;
\item[(b)] $L$ est engendré par les coordonnées, dans une base de $X$ sur $K$, des vecteurs horizontaux de $X \otimes_K L$, et
\item[(c)] le corps des constantes de $L$ est $k$.
\end{enumerate}

La propriété de $(V, \nabla)$ : ``$V \otimes_K L$ est engendré sur $L$ par ses vecteurs horizontaux'' est stable par produits tensoriels et sous-quotients. Si elle est vérifiée pour $V$, on a
\[
(V \otimes_K L)^\nabla \otimes_k L \xrightarrow{\sim} V \otimes_K L
\]
et en particulier, si $k$ est le corps des constantes de $L$,
\[
(V \otimes_K L)^\nabla \otimes_k L \xrightarrow{\sim} V \otimes_K L.
\]

Si $(L, \partial)$ est une extension de Picard-Vessiot de $K$, pour $(X, \nabla)$, le foncteur
\[
\omega_L : \langle X \rangle^\otimes \to \Vect(k) : V \mapsto (V \otimes_K L)^\nabla
\]
est donc un foncteur fibre.

Sur $L$, on dispose de l'isomorphisme (9.4.1) entre les foncteurs fibres déduit de $\omega_0$ et $\omega_L$: un $K$-morphisme
\begin{equation}\label{eq:9.4.2}
\Spec(L) \to \underline{\Isom}_k(\omega_L \otimes K, \omega | \langle X \rangle^\otimes).
\end{equation}

\begin{proposition}[9.5]
Avec les notations précédentes, (9.4.1) identifie $L$ au corps des fonctions rationnelles de $\underline{\Isom}_k(\omega_L \otimes K, \omega | \langle X \rangle^\otimes)$.
\end{proposition}

\begin{proof}
Soient $G = \underline{\Aut}_k(\omega_0)$ et $P$ le $G_K$-torseur $\underline{\Isom}(\omega_L \otimes K, \omega | \langle X \rangle^\otimes)$. Le morphisme (9.4.2), en termes d'anneaux, est
\begin{equation}\label{eq:9.5.1}
\Gamma(P, \mathcal{O}) \to L.
\end{equation}
Le noyau $I$ de ce morphisme est un idéal $\nabla$-stable. Par (9.3.2), $I = 0$ et (9.5.1) est injectif. Le $G$-torseur $P$ est contenu dans le $\GL(\omega_L(X))$-torseur $\underline{\Isom}(\omega_L \otimes K, X)$ et que $\Gamma(P, \mathcal{O})$ engendre $L$ exprime 9.4(b).
\end{proof}

\textbf{Scholie (9.6).} Si $\omega_0 : \langle X \rangle^\otimes \to \Vect(k)$ est un foncteur fibre, le torseur $P$ de 9.5 a pour corps de fonctions rationnelles $k(P)$ une extension de Picard-Vessiot de $K$: pour $V$ dans $\langle X \rangle^\otimes$ et $a \in \omega_0(X)$, $p(a)$ est une section horizontale de $V$ sur $P$, d'où
\begin{equation}\label{eq:9.6.1}
a \mapsto p(a) : \omega_0(V) \xrightarrow{\sim} (V \otimes_K L)^\nabla
\end{equation}
\begin{equation}\label{eq:9.6.2}
\omega_k(V) \otimes_k L \xrightarrow{\sim} V \otimes_K L.
\end{equation}
Pour $V = X$, ceci fournit (a), l'assertion (b) résulte de ce que $P \subset \underline{\Isom}_k(\omega_0(X), X \otimes K)$ et (c) est 9.3(ii). Par (9.6.1), $\omega_0$ est canoniquement isomorphe à $\omega_L$ (6.4). Ceci, et 9.4, donnent un dictionnaire entre extensions de Picard-Vessiot pour $(X, \nabla)$ et foncteurs fibres pour $\langle X \rangle^\otimes$.

\begin{corollary}[9.7]
Si $k$ est algébriquement clos, il existe une extension de Picard-Vessiot $L$ pour $(X, \nabla)$. Elle est unique à isomorphisme (non unique) près.
\end{corollary}

Soit $G := \underline{\Aut}_k(\omega_L)$. Le groupe des $K$-automorphismes différentiels de $L$ est $G(k)$: $G$ est le \emph{groupe de Picard-Vessiot}. Voir I.~Kaplansky, \textit{An introduction to differential algebra}, Hermann, Paris 1957.

\begin{remark}[9.8]
Soient $(e^i)_{1 \leq i \leq n}$ une base de $X$ et $\nabla e^i = \sum_j c^i_j e^j$. Soit $M$ une extension différentielle de $K$. Pour que $\xi = \sum y_i \otimes e^i \in X \otimes_K M$ soit horizontal, il faut et suffit que les $y_i$ vérifient le système
\begin{equation}\label{eq:9.8.1}
\partial y_i = -\sum_j y_j c^j_i.
\end{equation}
Plutôt que d'une extension de Picard-Vessiot pour $X$, on parle classiquement d'une extension de Picard-Vessiot pour un système (9.8.1): une extension différentielle de $L$, engendrée par des $y_{\alpha i}$ formant $n$ solutions linéairement indépendentes de (9.8.1) et de même corps de constantes que $K$.

Plus classiquement encore, on considère des équations différentielles linéaires du $n$\textsuperscript{ème} ordre en une indéterminée $y$. On passe de là à un système par l'astuce habituelle: passer au système vérifié par $(y, \partial y, \partial^2 y, \ldots, \partial^{n-1}y)$.
\end{remark}

\begin{remark}[9.9]
Plus généralement, soient $K$ un corps de caractéristique 0, $\mathcal{D}$ un espace vectoriel sur $K$ et $[,]$ un crochet de Lie sur $\mathcal{D}$. On ne le suppose pas $K$-linéaire, mais on suppose qu'il existe $\delta : \mathcal{D} \to \Der(K, K)$ (linéaire et morphisme de Lie) vérifiant
\begin{equation}\label{eq:9.9.1}
[D_1, \lambda D_2] - \lambda[D_1, D_2] = \delta(D_1)(\lambda) D_2.
\end{equation}
Le corps des constantes $k$ est l'intersection des noyaux des $\delta(D)$ ($D \in \mathcal{D}$).

Une \emph{connexion intégrable} sur un $K$-espace vectoriel $V$ est une action de Lie $\nabla$ de $\mathcal{D}$ sur $V$ vérifiant
\[
\nabla_D(\lambda v) = \delta(D)(\lambda) \cdot v + \lambda \nabla_D v.
\]
Les $K$-espaces vectoriels de dimension finie à connexion intégrable $(V, \nabla)$ formant une catégorie tannakienne sur $k$.

Une \emph{extension différentielle} de $(K, \mathcal{D})$ est $(L, \mathcal{D}')$ avec $L$ extension de $K$, muni de $L \otimes_K \mathcal{D} \xrightarrow{\sim} \mathcal{D}'$ compatible aux actions $\delta$ de $\mathcal{D}$ et $\mathcal{D}'$ sur $K$ et $L$.

Une \emph{extension de Picard-Vessiot} de $(K, \mathcal{D})$, pour $(V, \nabla)$, est une extension différentielle $(L, \mathcal{D}')$ de $(K, \mathcal{D})$, de même corps de constantes, telle que $L \otimes_K V$ soit engendré par ses vecteurs horizontaux (annulés par les $\nabla_D$) et que $L$ soit engendré par les coordonnées, dans une base de $V$ sur $K$, des vecteurs horizontaux de $V \otimes_K L$. Comme en 9.6, les extensions de Picard-Vessiot correspondent aux foncteurs fibres de $\langle V \rangle^\otimes$ sur $k$. En particulier, si $k$ est algébriquement clos, elles existent et sont uniques à isomorphisme près.

\textit{Exemple:} on prend pour $K$ un corps aux dérivées partielles: muni de $d$ dérivations $\partial_i$ qui commutent. On prend pour $\mathcal{D}$ le $K$-espace vectoriel de base les $\partial_i$, pour $\delta : \partial_i \mapsto \partial_i$ et on définit $[,]$ par (9.9.1) et $[\partial_i, \partial_j] = 0$.
\end{remark}

\section*{Index terminologique}

\begin{tabbing}
\hspace{3cm} \= \hspace{3cm} \= \kill
anneau (algèbre): à unité \> \> 1.1 \\[4pt]
Barr-Beck (théorème de) \> \> 4.1 \\[4pt]
coaction (d'un cogébroïde) \> \> 1.15 \\
\> (d'une comonade) \> 4.1 \\[4pt]
cogébroïde \> \> 1.15 \\
\> (des $k$-endomorphismes de $\omega$) \> 4.7 \\[4pt]
comonade \> \> 4.1 \\[4pt]
compact closed \> \> 2.5 \\[4pt]
contragrédient \> \> 2.4 \\[4pt]
contrainte d'associativité (commutativité) \> \> 1.2 \\[4pt]
\> d'unité \> 2.1.1 \\[4pt]
dimension \> \> 7.1 \\[4pt]
dual \> \> 2.2 \\[4pt]
fibre (foncteur --): $\otimes$-foncteur exact $k$-linéaire \> \> 1.9 \\
\> fidèlement plat (dans $\mathcal{T}$) \> 7.7/7.8 \\[4pt]
$\otimes$-foncteur; morphisme de --- \> \> 2.7 \\[4pt]
fppf: fidèlement plat de présentation finie \> \> \\
fqc: fidèlement plat quasi-compact \> \> \\[4pt]
gerbe \> \> 3.2 \\
gerbe à lien affine \> \> 3.6 \\
$\mathcal{T}$-groupe \> \> 7.8 \\
($k$-)groupoïde \> \> 1.6 \\[4pt]
Hom interne \> \> 2.2 \\[4pt]
Induit (groupoïde ---) \> \> 1.6 \\
Ind-objet \> \> 7.5 \\[4pt]
$\otimes$-isomorphisme (isomorphisme de foncteur fibres) \> \> 1.9, cf.~2.7 \\
localement (dans $\mathcal{T}$): pour la topologie fqc \> \> 7.8 \\
module (sur un $\mathcal{T}$-schéma) \> \> 7.8 \\
monoïdal (catégorie ---) \> \> 2.2 \\[4pt]
puissance extérieure, symétrique \> \> 7.1, 7.8 \\
\> (dans une catégorie tensorielle) \> 7.7 \\
(fidèlement) plat (dans une catégorie tensorielle) \> \> 1.15 \\
représentation (d'un cogébroïde) \> \> 3.2 \\
\> (d'une catégorie fibrée, d'une gerbe) \> 3.6 \\
\> (d'un groupoïde) \> 1.6 \\
rigide \> \> 2.5 \\
schéma affine en $\mathcal{T}$ = $\mathcal{T}$-schéma affine \> \> 7.8 \\[4pt]
sections globales d'un module sur un $\mathcal{T}$-schéma \> \> 7.8 \\
$\mathcal{T}$-schéma vectoriel \> \> 7.10 \\
super- \> \> 1.4 \\
tannakien (catégorie ---) \> \> 2.8 \\
tensoriel (catégorie ---) \> \> 1.2, 2.1 \\
\> (produit --- de catégories abéliennes) \> 5.1 \\
transitif (groupoïde ---) \> \> 1.6 \\
transposé \> \> 2.4 \\
vide ($\mathcal{T}$-schéma --- non) \> \> 7.8
\end{tabbing}

\section*{Index des notations}

Pour l'usage de $\Isom$, $\Aut$, avec $\otimes$ en exposant et $k$ ou $S$ en indice, voir 1.11. Pour $\Hom$ et $\End$: 6.5.

\begin{tabbing}
\hspace{3cm} \= \kill
ACU: \> 2.7 \\[4pt]
$(A \to B)$, $(B \to A)$ \> 5.10 \\
$\coh$ : $(A)_{\coh}$ \> 5.2 \\[4pt]
$\eed$ : $\Hom^{\eed}$ : catégorie de foncteurs exacts à droite \> 5.2 \\
$\otimes_k$ : $\otimes$ \> 5.2 \\[4pt]
$\Ind(\cdots)$: catégorie des ind-objets (cf.~7.5) de $\cdots$ \>  \\
$k$: toujours commutatif \\[4pt]
$L_k(\omega_1, \omega_2)$ , $L(\omega)$ \> 4.7 \\
$\Lambda(\omega_1, \omega_2)$, $\Lambda(\omega)$ (dans $\mathcal{T}$) \> 8.4 \\
``$\varinjlim$'' \> 7.5 \\[4pt]
$(\mathrm{Pt})$ (dans $\mathcal{T}$) \> 7.8 \\
$\ptf$ : $(B\text{-}\mod)_{\ptf}$: projectif de type fini \> 6.1 \\[4pt]
$\Sp(\cdots)$ \> 7.8 \\
$\Rep(S : G)$ \> 3.2 \\
$\Rep(G)$ \> 1.5 \\[4pt]
$\mathcal{T}$-schéma affine, groupe \> 7.8 \\
schéma vectoriel \> 7.10 \\
$\Vect(k)$ \> 1.3 \\[4pt]
$\langle X \rangle$ \> 5.12 \\
$\langle X \rangle^\otimes$ \> 6.16 \\
$\bigwedge^n_\mathcal{T}$ \> 7.1
\end{tabbing}

\section*{Bibliographie}

\begin{enumerate}
\item[{[1]}] M.~Artin, \textit{Versal deformations and algebraic stacks}, Inv.~math.~27 (1974) pp.~165--189.

\item[{[2]}] P.~Deligne and J.S.~Milne, \textit{Tannakian categories} in Hodge Cycles, Motives, and Shimura Varieties, Lecture Notes in Mathematics 900, Springer Verlag, 1982, pp.~101--228.

\item[{[3]}] E.R.~Kolchin, \textit{Differential algebra and algebraic groups}, Academic Press, 1973.

\item[{[4]}] S.~MacLane, \textit{Natural associativity and commutativity}, Rice University Studies 49 (1963) pp.~28--46.

\item[{[5]}] S.~MacLane, \textit{Categories for the working mathematician}, Graduate texts in math.~5 (Springer-Verlag 1971).

\item[{[6]}] S.~Saavedra, \textit{Catégories tannakiennes}, Lecture Notes in Mathematics 265, Springer Verlag 1972.
\end{enumerate}

\subsection*{Sigles}

SGA : Séminaire de géométrie algébrique du Bois-Marie.

SGA 1: Revêtements étales et groupe fondamental: Lecture Notes in Mathematics 224, Springer Verlag, 1971.

SGA 3: Schémas en groupes, vol.~1: Lecture Notes in Mathematics 151, Springer Verlag, 1970.

SGA 4: Théorie des topos et cohomologie étale des schémas, vol.~1: Lecture Notes in Mathematics 269, Springer Verlag, 1972.

\vspace{1cm}

Juin 1990. Une approche toute différente de résultats proches de ceux du paragraphe 7 a été développée par S.~Doplicher et J.E.~Roberts (Ann.\ of Math.\ 130 p.~75--119 (1989) et Inv.\ Math.\ 98 p.~175--218 (1989)). Dans un langage un peu différent du leur: ils considèrent une catégorie tensorielle sur $\mathbb{R}$, $\mathcal{T}$, polarisée au sens de [6], et prouvent qu'elle est la catégorie des représentations d'un groupe compact, munie de sa polarisation naturelle. Le début de leur preuve, parallèle au début du paragraphe 7, observe que la dimension de chaque objet est un entier $\geq 0$ et que si $V_1, \ldots, V_k$ sont de dimension $n_1, \ldots, n_k$, il existe un $\otimes$-foncteur de la catégorie des représentations de $\prod \GL(n_i)$ dans $\mathcal{T}$ envoyant la représentation standard de $\GL(n_i)$ sur $V_i$. Le premier point acquis, leur résultat pourrait être déduit de ceux du paragraphe 7 et de Saavedra (Chap~VI). Leur preuve est toute différente.

\vspace{1cm}

\begin{flushright}
Institute for Advanced Study\\
Princeton, New Jersey 08540
\end{flushright}

\end{document}
